\documentclass[12pt,a4paper]{article}
\usepackage{../shared/pt-common}
\usepackage[margin=2.5cm]{geometry}
\usepackage{setspace}
\usepackage{tikz-3dplot}
\onehalfspacing

\title{\textbf{La Théorie de la Persistance~I~:\\Le Crible comme Système Dynamique}}
\author{Yan Senez\\\texttt{yan.senez@protonmail.com}}
\date{}

\begin{document}
\maketitle

\begin{abstract}
Ceci est le premier de cinq articles présentant les fondations mathématiques
de la Théorie de la Persistance (PT).  L'entrée unique est
le paramètre de symétrie $s = 1/2$, forcé par l'interdiction des
auto-transitions de la séquence de résidus mod-$3$ parmi les entiers
$6$-rugueux (Théorème~T1)~: la matrice de transfert
$T_3 = \mathrm{antidiag}(1,1)$ a une distribution stationnaire unique
en $(1/2, 1/2)$, sans paramètre libre.  Ce n'est pas une hypothèse~:
$s = 1/2$ est l'unique point fixe d'une matrice $2 \times 2$ doublement
stochastique antidiagonale (algèbre linéaire), et le théorème
d'équidistribution de Dirichlet confirme la convergence empirique.

À partir de ce point de départ, nous établissons quatre groupes de résultats.
\textbf{(i)}~Le crible comme champ dynamique (Section~\ref{sec:sieve})~:
la séquence de lacunes $\{g_n\}$ est l'unique degré de liberté~;
T1 force $T_3$ et $s = 1/2$ (distribution stationnaire unique de la
matrice antidiagonale).
\textbf{(ii)}~Unicité d'Ératosthène (Section~\ref{sec:uniqueness})~:
quatre théorèmes de naturalité (N1--N4) prouvent qu'Ératosthène est l'unique
crible auto-cohérent, minimal, primo-dynamique~; le Théorème~T6 fournit
trois preuves indépendantes---théorique des champs (T6a), axiomatique via quatre
axiomes de congruence d'anneau C1--C4 (T6b), et géométrique-informationnelle via
l'unicité de Shore--Johnson et Čencov (T6c).
\textbf{(iii)}~Lois de conservation (Section~\ref{sec:conservation})~:
la distribution de lacunes d'entropie maximale est géométrique (L0, par
concavité stricte)~; l'exposant de conservation spectrale satisfait
$\alpha_{\rm cons} = s^2 = 1/4$ (Théorème~T2)~; la bifurcation
sommet--arête produit exactement deux paramètres canoniques
$\qstat = 13/15$ et $\qtherm = e^{-1/15}$ au point fixe
$\mustar = 15$ (prouvé dans l'Article~III~\cite{companion_M3}).
\textbf{(iv)}~L'Identité Fondamentale des Lacunes (GFT) (Section~\ref{sec:gft})~:
l'identité algébrique (valide pour toute distribution sur $m$ classes)
$\log_2 m = \DKL + H$ décompose la
capacité informationnelle totale en structure persistante et bruit
irréductible~; la Flèche du Temps $dH/d\DKL = -1$ est une identité, non une
inégalité~; l'identification de Ruelle connecte la GFT au
principe variationnel du formalisme thermodynamique avec l'énergie libre
$F_R = 0$.

Tous les résultats utilisent la combinatoire finie, l'algèbre linéaire et le
théorème de Perron--Frobenius. Aucune physique n'est invoquée~; aucun paramètre n'est
ajusté. Les quatre articles suivants développent la géométrie informationnelle
et l'holonomie (M2), la convergence et le point fixe (M3), les structures
mathématiques étendues (M4), et le pont de l'arithmétique à la
physique (M5).
\end{abstract}

\tableofcontents
\newpage

% ============================================================
%  Introduction
% ============================================================
\section{Introduction}
\label{sec:intro}

Cet article présente la première couche des fondations mathématiques
de la Théorie de la Persistance (PT), un cadre qui étudie la structure
information-théorique et dynamique de la séquence de lacunes premières
à travers le crible d'Ératosthène.  L'observation centrale est que le
crible impose des contraintes exactes sur les transitions de résidus---plus
fondamentalement, l'interdiction absolue des auto-transitions modulo~3
(Théorème~T1)---à partir de laquelle une structure mathématique riche se déploie
sans paramètre libre.

Le développement est organisé comme une série de cinq articles~:
\begin{description}[nosep]
\item[M1 (cet article)~:] Le crible comme champ dynamique, unicité
  d'Ératosthène, lois de conservation, et le Théorème Fondamental des Lacunes.
\item[M2~:] Géométrie informationnelle, holonomie, et l'échelle d'angle interne.
\item[M3~:] Convergence, récurrences exactes, et le point fixe
  $\mustar = 15$.
\item[M4~:] Structures mathématiques étendues---algèbre du crible, mécanique complexe, cadre catégorique, mathématiques prédictives.
\item[M5~:] Le pont de l'arithmétique à la physique---dualité de Pontryagin, contraction de Buchstab, reconstruction d'Osterwalder--Schrader.
\end{description}

L'article présent établit le noyau arithmétique. Nous commençons par
identifier la séquence de lacunes $\{g_n = p_{n+1} - p_n\}$ comme l'unique
champ dynamique (Section~\ref{sec:sieve}).  La loi de
transition interdite T1 force la matrice de transfert mod-$3$ à être purement antidiagonale,
avec distribution stationnaire unique $s = 1/2$. Trois routes indépendantes vers
$T_3$ sont données (arithmétique des lacunes, involution $\mathbb{Z}/6\mathbb{Z}$,
contrainte spectrale).

La Section~\ref{sec:uniqueness} prouve que le crible d'Ératosthène est
le \emph{unique} crible compatible avec la structure multiplicative
de~$\mathbb{N}$. Quatre théorèmes de force croissante (N1--N4)
établissent l'unicité algébrique, l'auto-cohérence, la minimalité structurelle,
et l'ordre canonique qui force $p = 3$ comme premier niveau en cascade. Le Théorème~T6 fournit ensuite trois preuves indépendantes de l'unicité du crible~: de la structure de champ (T6a), des quatre
axiomes de congruence d'anneau (T6b), et de la géométrie informationnelle canonique (T6c).

La Section~\ref{sec:conservation} dérive les lois de conservation. La
distribution de lacunes d'entropie maximale unique est géométrique (L0)~; l'exposant de
conservation spectrale $\alpha_{\rm cons} = s^2 = 1/4$ contrôle
les taux de convergence (T2)~; la bifurcation sommet--arête en $\mustar = 15$
(prouvée dans l'Article~III~\cite{companion_M3})
produit exactement deux paramètres de déformation canoniques.

La Section~\ref{sec:gft} établit l'Identité Fondamentale des Lacunes (GFT)~:
$\log_2 m = \DKL(P\|U_m) + H(P)$, la décomposition de Shannon appliquée aux
distributions de résidus premiers. Le contenu mathématique est une
identité algébrique valide pour \emph{toute} distribution de probabilité sur $m$ classes~;
la nouveauté est l'\emph{application} à la dynamique du crible, pas l'algèbre. L'équivalence de Ruelle, la Flèche du Temps, et la limite de Bekenstein
suivent comme corollaires immédiats.

Toutes les preuves utilisent seulement la combinatoire finie, l'algèbre linéaire et le
théorème de Perron--Frobenius. Aucune physique n'est invoquée nulle part. Les
théorèmes numérotés sont~: T1 (transitions interdites), T3 (transfert
antidiagonal), T6 (unicité du crible, trois preuves), L0 (entropie maximale),
T2 (conservation spectrale), GFT (identité informationnelle), et la
chaîne de naturalité N1--N4.


% ============================================================
%  S2. The Sieve as a Dynamical Field
% ============================================================
\section{Le Crible comme Champ Dynamique}
\label{sec:sieve}

\epigraph{God made the integers; all else is the work of man.}{Leopold Kronecker}

% ---- Status Box (R1) ----------------------------------------
\begin{statusbox}
\textbf{GOAL:}
  Establish the gap sequence $\{g_n\}$ as the unique dynamical field
  of the Theory of Persistence and derive the exact mod-3 transition law T1.

\textbf{INPUTS:}
  None. This is the entry point.

\textbf{MAIN RESULT:}
  T1 (Theorem~\ref{thm:T1}): self-transitions mod~3 are absolutely forbidden.
  $T_3 = \mathrm{antidiag}(1,1)$ (Theorem~\ref{thm:T3}).
  $s = 1/2$ (unique stationary distribution of $T_3$; Theorem~\ref{thm:T3}).

\textbf{STATUS:}
  \THMTAG~(T1, T3, $s = 1/2$, Unique DOF, T1-3gram)

\textbf{NOT PROVED:}
  Convergence $\alpha_k \to 1/2$ (that is T5, the companion article~\cite{companion_M3}).
  Uniqueness of the Eratosthenes sieve (\cref{sec:uniqueness}).
  Any physical claim.
\end{statusbox}

\begin{remark}[Ce que PT ne \emph{pas} affirme]
\label{rem:ch1_limits}
Avant de procéder, nous déclarons explicitement ce qui se trouve \emph{en dehors}
du champ de ce travail~:
\begin{enumerate}[nosep]
\item PT n'affirme pas être une « Théorie de Tout ».
  Le pont vers les observables physiques (43 quantités du Modèle Standard à
  $0,30\%$ d'écart moyen, zéro paramètres ajustés) est développé dans
  l'Article~V~\cite{companion_M5}~; cet article dérive zéro
  observables physiques. PT ne dit rien sur la conscience, l'éthique ou
  l'expérience subjective.
\item PT ne prouve pas que l'univers physique \emph{est} le crible.
  Il prouve que la structure du crible correspond aux constantes du MS.
  Si cette correspondance est accidentelle, structurelle ou ontologique est
  une question philosophique ouverte.
\item PT ne remplace pas la théorie quantique des champs. Il reconstruit
  le cadre QFT (espace de Hilbert, règle de Born, règles de Feynman) à partir
  des données du crible, mais utilise le langage de QFT pour exprimer les prédictions physiques.
\item PT does not explain \emph{why} mathematics exists.  The mystery
  recedes one level: from ``why these constants?''\ to ``why these
  mathematical structures?''
\item Chaque résultat porte une étiquette épistémique explicite
  (\THMTAG, \IDTAG, \DERTAG, \BRIDGETAG, etc.).
  Les affirmations étiquetées \BRIDGETAG\ sont des identifications structurelles, non
  des théorèmes~; les affirmations étiquetées \CONDTAG\ dépendent des hypothèses énoncées.
  Le lecteur doit toujours vérifier l'étiquette avant d'évaluer une affirmation.
\end{enumerate}
\end{remark}

\begin{remark}[Convergence indépendante]
\label{rem:ch1_convergence}
Le pont entre la théorie des nombres et la physique fondamentale que PT
constitue systématiquement a récemment attiré un soutien indépendant
de plusieurs directions~:
\begin{itemize}[nosep]
\item \textbf{L'annulation d'anomalies est arithmétique.}
  Camp, Gripaios and Nguyen~\cite{CampGripaiosNguyen2024} proved that
  abelian gauge anomaly cancellation in 2D admits solutions if and only
  if it admits solutions in $\mathbb{Q}_p$ for every prime~$p$
  (Hasse--Minkowski principle).
  Lee, Takahashi and Tsai~\cite{LeeTakahashiTsai2026} showed that
  anomaly cancellation for minicharged particles is equivalent to the
  Prouhet--Tarry--Escott problem at degree~3.
  In both cases, quantum consistency reduces to a number-theoretic
  constraint---exactly as PT derives it from the sieve.
\item \textbf{Les zéros de zêta sont des spectres physiques.}
  Remmen~\cite{Remmen2021} constructed a scattering amplitude whose
  mass tower coincides with the Riemann zeta zeros, with the critical
  line $\Re(s) = 1/2$ playing the role of a unitarity bound.
\end{itemize}
Ces résultats ont été obtenus sans connaissance de PT et par des
métres entièrement différentes. Ils ne prouvent pas PT, mais ils établissent
que sa prémisse centrale---la théorie des nombres dicte la physique---
gagne une traction indépendante.
\end{remark}

% ============================================================
\subsection{Intuition~: le crible comme système dynamique}
\label{sec:ch1_intuition}

Le Crible d'Ératosthène est généralement présenté comme un algorithme pour
générer les nombres premiers~: rayer les multiples de~2, puis~3, puis~5, et ainsi
de suite. Ce qui reste, ce sont les nombres premiers. C'est correct mais incomplet. Le
crible ne fait pas que \emph{filtrer}---il \emph{contraint} aussi. À
chaque étape, l'élimination des composites impose des règles de sélection sur les
classes de résidus qui survivent à chaque étape ultérieure. Ces règles de
sélection s'accumulent en un système dynamique sur la séquence de lacunes
$g_n = p_{n+1} - p_n$ entre les nombres premiers consécutifs.

\begin{remark}[Zéro est une certitude, non le néant]
\label{rem:zero_certainty}
Le crible sépare \textbf{la certitude de l'incertitude}.

Un composé $n = kp$ satisfait $n \equiv 0 \pmod{p}$~: son
cycle $p$ s'est \emph{fermé}---revenu à l'origine après $k$ tours
complets de $\mathbb{Z}/p\mathbb{Z}$. Nous \emph{savons} qu'il est divisible
par~$p$~; il n'y a aucune incertitude sur son résidu. Le résidu zéro
n'est pas l'absence mais \textbf{certitude}~: l'état d'annulation de phase
complète, où le cycle est terminé et la réponse
est déterminée.

Un nombre premier $p$ est l'opposé~: aucun cycle ne se ferme en dessous de~$p$. Ses résidus
modulo tous les nombres premiers plus petits sont non-nuls---\emph{phases
ouvertes} perpétuelles, portant une incertitude non résolue.
Aucun test modulaire ne prédit un nombre premier~; il est \emph{irréductiblement incertain}.

La hiérarchie de certitude est donc~:
\begin{center}
\begin{tabular}{cl}
\toprule
\textbf{Objet} & \textbf{Statut} \\
\midrule
$0$     & Certitude absolue (chaque cycle fermé, $H = 0$) \\
$1$     & Neutre (aucun nombre premier ne divise~; aucune information) \\
Composés & Certitude partielle (quelques cycles fermés) \\
Nombres premiers & Incertitude maximale (aucun cycle ne se ferme) \\
\bottomrule
\end{tabular}
\end{center}
Le crible \emph{supprime les certitudes} (cycles fermés) et
\emph{conserve les incertitudes} (phases ouvertes). Les nombres premiers survivent
précisément parce qu'ils sont \emph{imprévisibles} depuis en bas.

Cette inversion---zéro comme certitude, nombres premiers comme incertitude---est
fondamentale. Elle réapparaît comme~:
\begin{itemize}[nosep]
\item L'identité GFT (\cref{sec:gft})~: $\DKL + H = \log_2 m$,
  où $\DKL$ (structure = certitude) et $H$ (entropie = incertitude)
  sont des duaux conservés.
\item L'opérateur binaire $p = 2$ (\cref{rem:p2_operator})~:
  la première certitude (pair/impair est \emph{connu}), qui crée
  la partition info/anti-info.
\item L'opérateur d'incertitude $p = 3$~: la première \emph{incertitude}
  ($T_3$ antidiagonale, $s = 1/2$, alternation imprévisible),
  qui crée la première dimension spatiale.
\item L'équation de point fixe (l'article compagnon~\cite{companion_M3})~:
  une \emph{fermeture de cycle global}---le crible se reconnaît lui-même.
\end{itemize}
\end{remark}

\begin{remark}[Genèse~: de deux certitudes à toute la physique]
\label{rem:genesis}
Les entiers $0$ et $1$ sont deux \textbf{pôles distincts de certitude}~:
\begin{itemize}[nosep]
\item $0$~: certitude \emph{locale}---chaque nombre premier divise~$0$~;
  chaque cycle est fermé~; chaque question répondue.
\item $1$~: certitude \emph{globale}---aucun nombre premier ne divise~$1$~;
  c'est l'identité, l'élément neutre, certain de ne rien être
  en particulier.
\end{itemize}

Entre ces deux pôles, il y a exactement \textbf{deux états}~:
divisible ou non, fermé ou ouvert, zéro ou un. Ce compte---$2$---n'est
pas un postulat~; c'est le \emph{nombre de positions entre les deux
certitudes}. C'est le premier nombre premier $p_1 = 2$~: l'opérateur binaire.

Avoir deux états qui ne sont pas identiques force une
\textbf{troisième structure}~: la transition entre eux. Deux états
$A, B$ qui alternent créent une séquence $A B A B \ldots$ dont les
résidus modulo~$3$ parcourent $\{1, 2, 1, 2, \ldots\}$, sans jamais
frapper la même classe deux fois de suite (Théorème~T1 ci-dessous).
Le nombre $3 = p_2$ est la \emph{conséquence arithmétique} d'avoir
$2$ états binaires. C'est l'opérateur d'incertitude~: le premier
niveau où le résultat n'est pas prédéterminé par la structure binaire.

La chaîne causale est donc~:
\begin{equation}
  \{0, 1\} \;\text{exist}
  \;\;\Longrightarrow\;\;
  2 = p_1
  \;\;\xrightarrow{\text{T1}}\;\;
  s = \tfrac{1}{2}
  \;\;\Longrightarrow\;\;
  3 = p_2
  \;\;\Longrightarrow\;\;
  5, 7 = p_3, p_4
  \;\;\Longrightarrow\;\;
  \mustar = 15
  \;\;\Longrightarrow\;\;
  \text{SM.}
  \label{eq:genesis}
\end{equation}
L'\emph{véritable} entrée de la Théorie de la Persistance n'est pas $s = 1/2$ mais l'existence
de deux certitudes distinctes. Tout le reste---les nombres premiers,
le crible, l'holonomie, les observables physiques (Article~V~\cite{companion_M5})---découle de l'arithmétique
de l'espace entre $0$ et~$1$.
\end{remark}

\begin{remark}[Deux directions, une conséquence]
\label{rem:two_directions}
L'espace entre $0$ et $1$ peut être parcouru en
\textbf{deux directions}~:
\[
  \underbrace{1 \to 0}_{\text{openness} \to \text{certainty}}
  \qquad\text{and}\qquad
  \underbrace{0 \to 1}_{\text{certainty} \to \text{openness}}
\]
Les deux voient les mêmes $2$ états entre les pôles~; les deux créent le
même $p_1 = 2$~; les deux mènent au même T1, à la même $s = 1/2$,
à la même physique. Mais ils le voient depuis des \emph{côtés opposés}.

These two directions \emph{are} the bifurcation
$\qrel \neq \qtherm$ (\cref{sec:conservation}):
\begin{align}
  q_{\rm rel} &= 1 - \frac{2}{\mu}
  &&\text{(\textbf{relativistic branch}: builds structure,}
  \notag\\
  &&&\text{discrete, carries the binary $p_1 = 2$)}\,, \label{eq:q_rel} \\[4pt]
  q_{\rm therm} &= e^{-1/\mu}
  &&\text{(\textbf{thermodynamic branch}: dissipates structure,}
  \notag\\
  &&&\text{continuous, Boltzmann limit)}\,. \label{eq:q_therm}
\end{align}

La branche relativiste porte le facteur~$2 = p_1$ dans sa
définition~: l'opérateur binaire crée le réseau $2\mathbb{N}$
sur lequel la cascade opère. Il gouverne les couplages
($\alphaEM$, PMNS), les amplitudes de vertex, et le secteur info.
La branche thermodynamique est la limite continue où l'espacement
du réseau s'annule~; elle gouverne la géométrie (métrique, CKM),
les propagateurs d'arête, et le secteur d'entropie.

Les deux branches produisent la \emph{même conséquence}---la même
$s = 1/2$, les mêmes nombres premiers, le même $\mustar = 15$---car
la symétrie $0 \leftrightarrow 1$ (échange des deux
certitudes) est exacte. Cette symétrie est \textbf{CPT}~:
\begin{itemize}[nosep]
\item \textbf{C} (conjugaison de charge)~: échangez le \emph{contenu}
  (info $\leftrightarrow$ anti-info).
\item \textbf{P} (parité)~: échangez les \emph{pôles}
  ($0 \leftrightarrow 1$).
\item \textbf{T} (renversement du temps)~: échangez les \emph{directions}
  ($0{\to}1 \leftrightarrow 1{\to}0$).
\end{itemize}
L'invariance CPT est exacte car la physique dépend uniquement de l'\emph{espace entre}
les deux certitudes, non de quelle certitude
est appelée $0$ et laquelle est appelée~$1$.

\medskip\noindent
\textbf{Terminologie.}
Nous renommons la traditionnelle $\qstat$ (« statistique », opaque) en
$q_{\rm rel}$ (« relativiste »), reflétant son rôle de branche
de construction de structure qui porte l'opérateur binaire et
crée l'espace-temps. Le nom $q_{\rm therm}$ (« thermodynamique »)
est conservé.

\medskip\noindent
\textbf{Note épistémique.}
Les identifications « relativiste » et « thermodynamique » sont
des \emph{affirmations de pont} (voir l'article compagnon~\cite{companion_M5}), non des
théorèmes arithmétiques. Le contenu arithmétique de cette remarque est seulement que
deux directions existent entre $0$ et $1$, produisant deux valeurs
distinctes de~$q$. Les étiquettes physiques anticipent l'identification du pont
développée dans l'article compagnon~\cite{companion_M5}.
\end{remark}

\begin{remark}[Le crible comme dualité addition/multiplication]
\label{rem:add_mult_duality}
Le Crible d'Ératosthène alterne entre deux opérations
à chaque niveau~$p$~:
\begin{itemize}[nosep]
\item \textbf{Multiplication} ($\times p$)~: supprimez tous les multiples de~$p$.
  Cela crée un motif \emph{périodique} avec la période~$p$---un
  \textbf{cycle} (la roue au niveau~$p$).
\item \textbf{Addition} ($+1$)~: avancez vers l'entier suivant et
  mesurez la lacune jusqu'au prochain survivant.
\end{itemize}

L'interaction est constructive~:
\begin{enumerate}[nosep]
\item $\times p$ crée la \emph{boucle}~: la roue pour le
  primorial $p_1 \cdots p_k$ se répète avec la période $p_1 \cdots p_k$.
\item $+$ crée des \emph{nouvelles valeurs de lacune} à chaque niveau~:
  les nouvelles lacunes sont des \textbf{sommes} de lacunes précédentes
  ($4 = 2{+}2$, $6 = 4{+}2$, $10 = 4{+}6$).
\item La boucle s'\emph{épuise} à $p_{k+1}^2$~: au-delà de ce
  seuil, les composés de nombres premiers plus grands brisent la périodicité.
  Le nombre premier suivant la restaure avec une période plus longue.
\end{enumerate}

Numériquement, 59 des 60 premiers écarts de nombres premiers peuvent être exprimés comme
$g_n = g_i + g_j$ ou $g_n = g_i \times g_j$ d'écarts antérieurs.
La seule exception est $g_0 = 1$ (la graine). La multiplication
est toujours \emph{redondante} avec l'addition---chaque fois qu'un écart est un
produit, c'est aussi une somme---mais elle crée la périodicité que
l'addition ne peut pas.

Cette dualité fait écho à la Remarque~\ref{rem:zero_certainty}~:
$0$ est l'élément neutre de $+$, $1$ est l'élément neutre
de $\times$, et le crible alterne entre eux.
La bifurcation $\qrel \neq \qtherm$
(\cref{rem:two_directions}) est cette même dualité projetée~:
$\qrel = (\mustar{-}2)/\mustar$ (quotient, multiplicatif),
$\qtherm = e^{-1/\mustar}$ (série de Taylor, additive).
\end{remark}

L'idée clé est simple~: les nombres premiers consécutifs ne peuvent pas partager une classe
de résidu modulo~3. Si $p_n \equiv r \pmod{3}$ et $p_{n+1} \equiv r
\pmod{3}$, alors leur différence $g_n = p_{n+1} - p_n \equiv 0
\pmod{3}$, ce qui force $3 \mid g_n$. Pour les lacunes de taille au moins~2
et les nombres premiers supérieurs à~3, cela signifie que l'entier $p_n + g_n/2$
serait divisible par~3---contredisant sa primalité si $g_n \geq 6$,
ou menant à un point médian composé pour les petits multiples pairs de~3.
L'argument détaillé est le Théorème~T1 ci-dessous~; le processus de crible complet
est visualisé à la Figure~\ref{fig:sieve_process}.

La séquence de lacunes $\{g_n\}$ n'est pas simplement une quantité dérivée. C'est l'\emph{unique} degré de liberté indépendant~: donnés les lacunes et le
nombre premier initial $p_1 = 2$, la séquence de nombres premiers entière est reconstruite~;
conversement, connaître toutes les classes de résidu modulo chaque nombre premier récupère
les lacunes. Tout dans la Théorie de la Persistance---les matrices de transfert,
les angles d'holonomie, les constantes de couplage---est une fonctionnelle de
$\{g_n\}$.

\begin{remark}[PT n'est ni local ni global]
\label{rem:ch1_not_local}
Le Théorème du Reste Chinois (CRT) factorise chaque matrice de transfert
$T_m = T_{p_1} \otimes \cdots \otimes T_{p_k}$ en facteurs nombre premier par nombre premier. C'est un \emph{outil} puissant, mais cela ne fait pas de PT une
théorie « locale » où chaque nombre premier agit indépendamment. Plusieurs
structures sont intrinsèquement collectives~:
\begin{enumerate}[nosep]
\item Le théorème H $S_{K+1} \geq S_K$~: une contrainte de monotonie
  sur la distribution \emph{complète}.
\item Annihilation spectrale $r_2(0) = 0$ (l'article compagnon~\cite{companion_M3})~: une propriété de
  $T_3$ dans son ensemble, non décomposable en facteurs par nombre premier.
\item La métrique de Fisher (l'article compagnon~\cite{companion_M2})~: une géométrie riemannienne continue
  qui \emph{est} le crible, non une approximation de celui-ci.
\item Holonomie $\sin^2(\theta_p)$ (l'article compagnon~\cite{companion_M2})~: une phase accumulée
  autour de $\mathbb{Z}/p\mathbb{Z}$---cohérence globale, non une
  quantité ponctuelle.
\item Le lemme d'invariance automorphe (l'article compagnon~\cite{companion_M3})~: rigidité de
  $R_p = \{0\}$ sous $\mathrm{Aut}(\mathbb{Z}/p\mathbb{Z})$.
\end{enumerate}
PT dissout la distinction local/global tout comme elle dissout la
distinction discret/continu~: la structure de produit CRT et les
contraintes collectives coexistent comme descriptions complémentaires d'un
objet arithmétique unique.
\end{remark}

\begin{figure}[htbp]
\centering
\begin{tikzpicture}[>=stealth, xscale=0.60]
  % --- Row 0: All integers 2..23 ---
  \node[font=\scriptsize\bfseries, anchor=east] at (-1.4, 4.2) {Integers};
  \foreach \n [count=\i from 0] in {2,3,4,5,6,7,8,9,10,11,12,13,14,15,16,17,18,19,20,21,22,23} {
    \node[font=\scriptsize, minimum width=0.55cm] at (\i, 4.2) {\n};
  }

  % --- Row 1: Sieve — composites struck, primes circled ---
  \node[font=\scriptsize\bfseries, anchor=east] at (-1.4, 2.8) {Sieve};
  % Composites (struck out) — index = n-2
  \foreach \i/\n in {2/4, 4/6, 6/8, 7/9, 8/10, 10/12, 12/14, 13/15, 14/16, 16/18, 18/20, 19/21, 20/22} {
    \node[font=\scriptsize, text=gray!45] at (\i, 2.8) {\n};
    \draw[gray!40, thin] (\i-0.18, 2.65) -- (\i+0.18, 2.95);
  }
  % Primes (circled) — index = n-2
  \foreach \i/\n in {0/2, 1/3, 3/5, 5/7, 9/11, 11/13, 15/17, 17/19, 21/23} {
    \node[font=\scriptsize\bfseries, thmblue] (pr\n) at (\i, 2.8) {\n};
    \draw[thmblue, semithick] (\i, 2.8) circle (0.27cm);
  }

  % --- Row 2: Gaps (between primes, with arrows) ---
  \node[font=\scriptsize\bfseries, anchor=east] at (-1.4, 1.3) {Gap $g_n$};
  \foreach \xa/\xb/\g in {0/1/1, 1/3/2, 3/5/2, 5/9/4, 9/11/2, 11/15/4, 15/17/2, 17/21/4} {
    \pgfmathsetmacro{\xm}{(\xa+\xb)/2}
    \draw[condorange, <->, semithick] (\xa+0.32, 1.3) -- (\xb-0.32, 1.3);
    \node[font=\scriptsize\bfseries, condorange, fill=white, inner sep=1pt]
      at (\xm, 1.3) {\g};
  }

  % --- Row 3: Residues mod 3 ---
  \node[font=\scriptsize\bfseries, anchor=east] at (-1.4, -0.1) {$p \!\bmod 3$};
  \foreach \i/\r/\c in {0/2/predred, 1/0/valgray, 3/2/predred, 5/1/bridgepurple,
    9/2/predred, 11/1/bridgepurple, 15/2/predred, 17/1/bridgepurple, 21/2/predred} {
    \node[font=\scriptsize\bfseries, \c, draw=\c, rounded corners=2pt,
          minimum width=0.45cm, minimum height=0.38cm, inner sep=1pt]
      at (\i, -0.1) {\r};
  }

  % Annotation: alternation
  \draw[decorate, decoration={brace, amplitude=4pt, mirror}, thick]
    (3-0.35, -0.55) -- (21+0.35, -0.55);
  \node[font=\scriptsize, align=center] at (12, -1.0)
    {$p > 3$: residues alternate $1 \leftrightarrow 2$ (Theorem~T1)};

  % Vertical dashed guides from sieve row to residue row
  \foreach \i in {0, 1, 3, 5, 9, 11, 15, 17, 21} {
    \draw[gray!25, densely dotted] (\i, 2.42) -- (\i, 0.15);
  }
\end{tikzpicture}
\caption{Le processus du crible visualisé.
  \textbf{Rangée~1}~: entiers 2--23.
  \textbf{Rangée~2}~: composés barrés (Ératosthène), nombres premiers encerclés.
  \textbf{Rangée~3}~: lacunes $g_n = p_{n+1} - p_n$.
  \textbf{Rangée~4}~: résidus mod~3---après $p=3$, ils alternent strictement
  entre 1 et 2, interdits de se répéter (T1).}
\label{fig:sieve_process}
\end{figure}

% ============================================================
\subsection{La séquence de lacunes~: unique degré de liberté}
\label{sec:gap_sequence}

\begin{definition}[Séquence de Lacunes]
\label{def:gaps}
Soit $2 = p_1 < p_2 < p_3 < \cdots$ la séquence de tous les nombres premiers en
ordre croissant. La \emph{séquence de lacunes} est
\begin{equation}
  g_n = p_{n+1} - p_n, \qquad n \geq 1.
  \label{eq:gap_def}
\end{equation}
Les premières valeurs sont $g_1 = 1$, $g_2 = 2$, $g_3 = 2$, $g_4 = 4$,
$g_5 = 2$, $g_6 = 4$, $g_7 = 2$, $g_8 = 4$, $g_9 = 6$, \ldots
La lacune initiale $g_1 = 1$ (de $p_1=2$ à $p_2=3$) est exceptionnelle~;
toutes les lacunes ultérieures sont paires.
\end{definition}

\begin{definition}[Séquence de Résidu]
\label{def:residue}
Pour tout modulus entier $m \geq 2$, la \emph{séquence de résidu}
modulo~$m$ est
\begin{equation}
  r_n^{(m)} = p_n \bmod m, \qquad n \geq 1.
  \label{eq:residue_def}
\end{equation}
Quand $m$ est clair du contexte, nous écrivons $r_n$ pour $r_n^{(m)}$.
\end{definition}

La séquence de résidu est entièrement déterminée par la séquence de lacunes~:
\begin{equation}
  r_{n+1}^{(m)} = \bigl(r_n^{(m)} + g_n\bigr) \bmod m,
  \label{eq:gauge_connection}
\end{equation}
avec la condition initiale $r_1^{(m)} = 2 \bmod m$. Ce n'est pas une
approximation~; c'est une identité exacte.

\mTHM
\begin{theorem}[Unique Degré de Liberté]
\label{thm:unique_dof}
La séquence de lacunes $\{g_n\}_{n \geq 1}$ est l'unique champ dynamique
de \PT. Pour tout modulus $m$, toutes les séquences de résidu $\{r_n^{(m)}\}$,
toutes les matrices de transfert $T_m$, tous les angles d'holonomie $\theta_p$, et toutes
les constantes de couplage dérivées de ceux-ci sont des fonctionnelles mesurables
de~$\{g_n\}$.
\end{theorem}

\begin{proof}
La reconstruction de résidu~\eqref{eq:gauge_connection} est bijective~:
donnés $\{g_n\}$ et $r_1^{(m)} = 2 \bmod m$, la séquence entière
$\{r_n^{(m)}\}$ est déterminée. Conversement, $g_n = r_{n+1}^{(m)} -
r_n^{(m)} \bmod m$ récupère la lacune (modulo $m$). Puisque cela tient
pour tous les $m$ simultanément, la séquence de lacunes est le générateur
commun unique.

Pour la non-décomposabilité~: supposez $g_n = f_n + h_n$ avec $\{f_n\}$ et
$\{h_n\}$ statistiquement indépendants. Sous la carte de résidu mod-3, la
contrainte $P[r \to r] = 0$ (Théorème~\ref{thm:T1} ci-dessous) couple les paires consécutives
$(r_n, r_{n+1})$. Cet couplage est préservé par toute fonction de
$g_n$, donc tout sous-champ $\{f_n\}$ héritant de ce couplage ne peut pas
être indépendant de $\{h_n\}$. Le champ est indécomposable.

A complete formalization via the automorphic invariance lemma is given
in the companion article~\cite{companion_M3}: conditions U1--U4 uniquely determine the gap
sequence, subsuming both reconstruction and indecomposability.
\end{proof}

\begin{remark}
La forme de connexion de jauge~\eqref{eq:gauge_connection} n'est pas une
métaphore. C'est la propriété \emph{définissante} d'un champ de jauge de réseau valorisé en $\mathbb{Z}/m\mathbb{Z}$~: le résidu $r_n$ est le « transport parallèle »
de l'état initial $r_1$ le long du champ de lacune $\{g_n\}$, avec le
modulus $m$ sélectionnant le groupe de jauge. Nous rendons cela précis dans
l'article compagnon~\cite{companion_M2} (holonomie).
\end{remark}

% ============================================================
\subsection{Couverture additive des lacunes premières}
\label{sec:ch1_coverage}

La séquence de lacunes a une propriété algébrique remarquable~: elle est
\emph{additivement fermée}.

\begin{thmbox}{Génération Additive des Lacunes de Roue}
\label{thm:wheel_coverage}
Pour tout $k \geq 1$, l'ensemble des valeurs de lacune dans la roue au
niveau $k{+}1$ satisfait
$W_{k+1} \subseteq W_k \cup \mathrm{ADD}(W_k)$,
où $\mathrm{ADD}(G) = \{a + b : a, b \in G\}$.

\begin{proof}
En passant du niveau $k$ à $k{+}1$, le crible supprime
les multiples de $p_{k+1}$. Chaque suppression \emph{fusionne} deux lacunes
adjacentes $a, b \in W_k$ en une seule lacune $a + b$. Les lacunes dont
les extrémités sont toutes deux conservées passent inchangées. Par conséquent, chaque élément
de $W_{k+1}$ est soit dans $W_k$ soit dans $\mathrm{ADD}(W_k)$.
\end{proof}
\end{thmbox}

\begin{remark}[Extension aux lacunes premières (conjecture)]
\label{rem:coverage_conjecture}
Nous \emph{conjecturons} que la couverture additive s'étend aux lacunes
premières réelles~: pour tout $n \geq 2$,
$g_n \in G_{n-1} \cup \mathrm{ADD}(G_{n-1})$.
Ceci est vérifié informatiquement sur tous les $10^7$ écarts de nombres premiers en dessous
de $1,8 \times 10^8$ (couverture $100,0000\%$~; uniquement ADD $55,5\%$,
aussi générés multiplicativement $44,5\%$).
Le passage des lacunes de roue aux lacunes premières nécessite de contrôler
la primalité des extrémités, ce que le théorème de la roue n'aborde pas.

L'ensemble des valeurs de lacune distinctes croît comme $|G_N| \approx 0,13\,
(\ln p_N)^{2.26}$ et satisfait $\max(G_N) / |G_N| \approx 2$
across cinq ordres de grandeur.
La densité tronquée $\rho(G_N) = |G_N \cap [2, \max G_N]|
/ (\max G_N / 2) = 0,946$ pour $N = 10^7$ (105 des 111 valeurs
paires présentes~; aucun trou en dessous de~198).
Conditionnellement à la conjecture de Polignac (chaque entier pair est
une lacune première infiniment souvent), la densité de Schnirelmann de
l'ensemble infini $G_\infty = \bigcup_N G_N$ serait égale à~1,
prouvant la conjecture via le théorème de Schnirelmann.
Inconditionnellement, l'écart entre le théorème de la roue et la
conjecture complète reste un problème ouvert.
\end{remark}

% ============================================================
\subsection{Transitions interdites~: Théorème T1}
\label{sec:T1}

\mTHM
\begin{theorem}[Transitions Interdites (T1)]
\label{thm:T1}
Parmi les candidats du crible au niveau $p = 3$ (les entiers $6$-rugueux,
c.-à-d., les entiers copremiers aux deux~$2$ et~$3$), les auto-transitions de
la séquence de résidu mod-$3$ sont absolument interdites~:
\begin{equation}
  P\bigl[r_n \equiv 1 \to r_{n+1} \equiv 1 \pmod{3}\bigr] = 0, \qquad
  P\bigl[r_n \equiv 2 \to r_{n+1} \equiv 2 \pmod{3}\bigr] = 0.
  \label{eq:T1}
\end{equation}
Ceci vaut pour chaque paire d'entiers $6$-rugueux consécutifs, non seulement
asymptotiquement.
\end{theorem}

\begin{proof}
Les entiers $6$-rugueux sont exactement ceux de la forme $6k \pm 1$.
Leurs résidus modulo~3 alternent~: $6k - 1 \equiv 2$ et
$6k + 1 \equiv 1$. Les candidats consécutifs sont séparés par des lacunes
de 2 ou 4, et dans tous les cas, le résidu alterne entre 1 et 2.
Par conséquent, parmi les candidats du niveau-3 du crible, les auto-transitions sont
impossibles~: $P[r \to r] = 0$ exactement.
\end{proof}

\begin{corollary}[T1 pour les nombres premiers consécutifs]
\label{cor:T1_primes}
Comme chaque nombre premier $\geq 5$ est $6$-rugueux, T1 contraint les
transitions \emph{possibles} entre les nombres premiers consécutifs, mais n'interdit
pas tous les paires de même résidu~: les nombres premiers consécutifs peuvent partager un
résidu mod-$3$ (par exemple, $(61, 67)$, tous deux $\equiv 1 \pmod{3}$) car un
entier $6$-rugueux composé peut intervenir.
\end{corollary}

\begin{remark}[Utilisation PT de T1]
\label{rem:T0_scope}
Dans le cadre PT, T1 fixe la matrice de transition au niveau du crible~:
$T_3 = \mathrm{antidiag}(1,1)$, qui est exacte pour les entiers $6$-rugueux.
Le $T_3$ empirique pour les nombres premiers consécutifs a de petites entrées diagonales
non nulles, supprimées par $O(1/\ln N)$. Toutes les dérivations PT
(paramètre de symétrie $s = 1/2$, théorème de conservation, etc.)
reposent sur la matrice de niveau-crible, non sur la matrice empirique de niveau-premier.
L'écart par rapport à l'idéal diminue à mesure que $N$ croît (consistent avec
la conjecture d'Elliott--Halberstam, largement acceptée).
\end{remark}

% ============================================================
\subsection{La matrice de transfert exacte $T_3$}
\label{sec:T3}

\mTHM
\begin{theorem}[Transfert Antidiagonal (T3)]
\label{thm:T3}
Parmi les candidats du crible au niveau $p = 3$ (entiers $\equiv \pm 1 \pmod 6$),
la matrice de transition sur $\{1, 2\} \pmod 3$ est exactement
\begin{equation}
  T_3 = \begin{pmatrix} 0 & 1 \\ 1 & 0 \end{pmatrix}.
  \label{eq:T3}
\end{equation}
Le graphe de transition est montré à la Figure~\ref{fig:T3_graph}.
\end{theorem}

\begin{figure}[htbp]
\centering
\begin{tikzpicture}[>=stealth, scale=1.3]
  % Two state nodes
  \node[circle, draw=thmblue, thick, fill=thmblue!10,
        minimum size=1.2cm, font=\large] (s1) at (-1.8, 0) {$1$};
  \node[circle, draw=thmblue, thick, fill=thmblue!10,
        minimum size=1.2cm, font=\large] (s2) at (1.8, 0) {$2$};
  % Transition arrows (off-diagonal = 1)
  \draw[->, thick, thmblue] (s1) to[bend left=30]
    node[above, font=\scriptsize] {$P = 1$} (s2);
  \draw[->, thick, thmblue] (s2) to[bend left=30]
    node[below, font=\scriptsize] {$P = 1$} (s1);
  % Forbidden self-loops (crossed out)
  \draw[->, gray, dashed] (s1) to[out=140, in=220, looseness=5] (s1);
  \draw[predred, very thick] (-2.7, 0.9) -- (-2.3, 0.3);
  \draw[predred, very thick] (-2.7, 0.3) -- (-2.3, 0.9);
  \draw[->, gray, dashed] (s2) to[out=40, in=-40, looseness=5] (s2);
  \draw[predred, very thick] (2.3, 0.9) -- (2.7, 0.3);
  \draw[predred, very thick] (2.3, 0.3) -- (2.7, 0.9);
  % Matrix
  \node[font=\footnotesize, align=center] at (0, -1.8)
    {$T_3 = \begin{pmatrix} 0 & 1 \\ 1 & 0 \end{pmatrix}$
     \quad (antidiagonal)};
  % Labels
  \node[font=\scriptsize, gray] at (0, 1.2)
    {residue classes $\bmod\,3$};
  \node[font=\tiny, predred] at (-2.5, 0.0) {T1};
  \node[font=\tiny, predred] at (2.5, 0.0) {T1};
\end{tikzpicture}
\caption{Graphe de transition de la dynamique du crible mod-3.
  Sieve candidates alternate between residue classes $1$ and $2$
  with probability~1. Self-transitions are absolutely forbidden (T1,
  red crosses). The transfer matrix $T_3$ is purely antidiagonal---
  the simplest non-trivial dynamical law.}
\label{fig:T3_graph}
\end{figure}

\begin{proof}
Les candidats du crible au niveau~3 sont les entiers copremiers à $2 \times 3 = 6$,
c.-à-d., les entiers de la forme $6k \pm 1$ pour l'entier $k \geq 1$. Leurs résidus
modulo~3 sont~: $6k + 1 \equiv 1 \pmod 3$ et $6k - 1 \equiv 2 \pmod 3$.
Les candidats consécutifs de ce type sont $(6k-1, 6k+1)$ (lacune~2) ou $(6k+1, 6(k+1)-1) = (6k+1, 6k+5)$
(lacune~4). Dans les deux cas, les résidus alternent~: $1 \to 2$ ou $2 \to 1$. Aucune auto-transition
ne se produit jamais parmi les candidats du niveau-3 du crible. La stochasticité force les entrées hors-diagonales à~1.
\end{proof}

\begin{remark}[Trois routes indépendantes vers $T_3 = \mathrm{antidiag}(1,1)$]
\label{rem:T3_three_routes}
La forme antidiagonale de $T_3$ peut être établie par trois arguments disjoints,
utilisant chacun des primitives mathématiques différentes.

\smallskip\noindent
\textbf{Route 1~: Arithmétique des lacunes (ci-dessus).}
Énumérez les candidats $6k \pm 1$~; les seules lacunes possibles sont $2$ et $4$,
les deux échangent la classe de résidu modulo~3.

\smallskip\noindent
\textbf{Route 2~: Involution $\mathbb{Z}/6\mathbb{Z}$.}
Les survivants modulo~$6$ sont $\{1, 5\}$. L'application successeur $\sigma$ agit sur
cet ensemble à deux éléments et est donc l'unique involution non triviale~:
$\sigma(1) = 5,\; \sigma(5) = 1$. En réduisant modulo~3~:
$1 \mapsto 2,\; 2 \mapsto 1$, donc
$T_3 = \bigl(\begin{smallmatrix} 0 & 1 \\ 1 & 0 \end{smallmatrix}\bigr)$.
Cette preuve utilise seulement $|\mathrm{surv}_6| = 2$ et la structure d'involution,
sans calcul de lacune.

\smallskip\noindent
\textbf{Route 3~: Contrainte spectrale.}
$T_3$ est une matrice $2 \times 2$ doublement stochastique (par la symétrie
d'échange $1 \leftrightarrow 2$) avec $\mathrm{tr}(T_3) = 0$ (Théorème~T1~:
auto-transitions interdites). Les valeurs propres sont $\{+1, -1\}$.
La matrice doublement stochastique unique avec ces valeurs propres est
$\mathrm{antidiag}(1,1)$.
\end{remark}

\mCL
\begin{thmbox}{Paramètre de symétrie~: $s = 1/2$}
\label{thm:s_half}
La matrice antidiagonale $T_3 = \mathrm{antidiag}(1,1)$ (Théorème~T3)
a une distribution stationnaire unique~:
\begin{equation}
  s = \pi_1 = \pi_2 = \frac{1}{2}.
  \label{eq:s_half}
\end{equation}

\begin{proof}
L'équation du vecteur propre $\pi\,T_3 = \pi$ donne
$(\pi_1, \pi_2) \begin{psmallmatrix} 0 & 1 \\ 1 & 0 \end{psmallmatrix}
= (\pi_2, \pi_1) = (\pi_1, \pi_2)$,
donc $\pi_1 = \pi_2$. La normalisation $\pi_1 + \pi_2 = 1$ force
$s = 1/2$. Unicité~: la valeur propre $\lambda_0 = 1$ de la
matrice $2 \times 2$ doublement stochastique $T_3$ a la
multiplicité algébrique~1 (l'autre valeur propre est $\lambda_1 = -1$).
Par Perron--Frobenius, la distribution stationnaire est unique.
\end{proof}
\end{thmbox}

\begin{remark}[Épistémique de $s = 1/2$]
\label{rem:s_half}
La valeur $s = 1/2$ est un \textbf{théorème}, non une hypothèse. C'est
le point fixe unique de la matrice de transfert antidiagonale $T_3$,
prouvé ci-dessus par l'algèbre linéaire seule. Le théorème
d'équidistribution de Dirichlet confirme indépendamment que les fréquences
empiriques convergent vers cette valeur~: parmi les nombres premiers $> 3$, exactement la moitié
se trouvent dans chacune des deux classes de résidu non triviales modulo~3.
Le \emph{taux} de convergence ($\alpha_k \to 1/2$) est quantifié dans
l'Article~III~\cite{companion_M3} via l'annihilation spectrale.
\end{remark}

% ============================================================
\subsection{La connexion de jauge}
\label{sec:gauge}

L'équation~\eqref{eq:gauge_connection} se généralise aux moduli
sans carré arbitraires via le Théorème du Reste Chinois (CRT).

\begin{definition}[Modulus sans carré]
Un entier positif $m = q_1 \cdots q_k$ est \emph{sans carré} si
$q_1 < q_2 < \cdots < q_k$ sont des nombres premiers distincts.
\end{definition}

\mTHM
\begin{theorem}[Connexion de Jauge et Factorisation CRT]
\label{thm:gauge}
Soit $m = q_1 \cdots q_k$ sans carré. La récurrence de résidu
\[
  r_{n+1}^{(m)} = \bigl(r_n^{(m)} + g_n\bigr) \bmod m
\]
se décompose via l'isomorphisme $\mathbb{Z}/m\mathbb{Z} \cong
\prod_{i=1}^k \mathbb{Z}/q_i\mathbb{Z}$ (CRT) en $k$ composantes indépendantes,
chacune évoluant comme $r_{n+1}^{(q_i)} = (r_n^{(q_i)} + g_n)
\bmod q_i$. Les composantes à des nombres premiers distincts $q_i \neq q_j$ sont
statistiquement indépendantes (au premier ordre en $1/\ln N$).
\end{theorem}

\begin{proof}
L'isomorphisme CRT est un isomorphisme d'anneau, donc exact. Il mappe
$r \bmod m$ au vecteur $(r \bmod q_1, \ldots, r \bmod q_k)$.
La récurrence $r_{n+1} = r_n + g_n \pmod m$ mappe vers
$r_{n+1}^{(q_i)} = r_n^{(q_i)} + g_n \pmod{q_i}$ pour chaque $i$.
L'indépendance au premier ordre suit du théorème des nombres premiers
pour les progressions arithmétiques (PNT-AP)~: les résidus premiers modulo $q_i$ et
$q_j$ (pour les nombres premiers distincts $q_i, q_j$) sont asymptotiquement équidistribués
indépendamment. La correction est $O(1/\ln N)$ par le terme d'erreur PNT-AP.
\end{proof}

% ============================================================
\subsection{Monotonie de DPI}
\label{sec:dpi}

\mID
\begin{idbox}[Inégalité du Traitement des Données]
\label{id:dpi}
Si $m_1 \mid m_2$ (c.-à-d., $m_1$ divise $m_2$), alors
\begin{equation}
  \DKL(P_{m_1} \| U_{m_1}) \leq \DKL(P_{m_2} \| U_{m_2}),
  \label{eq:dpi}
\end{equation}
où $P_m$ est la distribution de résidu premier au modulus~$m$ et
$U_m$ est la distribution uniforme sur $(\mathbb{Z}/m\mathbb{Z})^\times$.

\begin{proof}
La projection naturelle $\pi_{m_2 \to m_1} : \mathbb{Z}/m_2\mathbb{Z}
\to \mathbb{Z}/m_1\mathbb{Z}$ (réduction modulo $m_1$) est un
homomorphisme d'anneau surjectif, donc un noyau de Markov. L'inégalité de
traitement des données pour la divergence KL énonce~: pour tout noyau
stochastique $K$, $\DKL(K(P) \| K(Q)) \leq \DKL(P \| Q)$. En appliquant à
$K = \pi_{m_2 \to m_1}$ avec $P = P_{m_2}$ et $Q = U_{m_2}$, et en notant que
$\pi_{m_2 \to m_1}(U_{m_2}) = U_{m_1}$ (l'uniforme mappe vers l'uniforme),
nous obtenons $\DKL(P_{m_1} \| U_{m_1}) \leq \DKL(P_{m_2} \| U_{m_2})$.
\end{proof}
\end{idbox}

La monotonie DPI dit~: plus de criblage = plus de structure = divergence plus élevée
par rapport à l'uniformité. Les niveaux du crible sont ordonnés par $\DKL$.

% ============================================================
\subsection{T1-3gram~: motifs interdits en trois étapes}
\label{sec:T1_3gram}

La loi de transition interdite T1 implique des contraintes sur les motifs
en trois étapes $(r_n, r_{n+1}, r_{n+2})$ de résidus mod-3.

\mTHM
\begin{theorem}[Validation Trigram (T1-3gram)]
\label{thm:T1_3gram}
Parmi les candidats du niveau-3 du crible, les motifs en trois étapes suivants sont
absolument interdits~:
\begin{equation}
  n_3(1,0,1) = 0, \qquad n_3(2,0,2) = 0.
  \label{eq:T1_3gram}
\end{equation}
Plus précisément~: des $27 = 3^3$ triplets mod-3 possibles, exactement
$12$ sont interdits~:
\begin{itemize}
\item $10$ par T1 (auto-transitions au sein du triplet), et
\item $2$ supplémentaires par la contrainte d'alternation de parité.
\end{itemize}
Les $15$ triplets restants sont autorisés. La symétrie de renversement du temps vaut
exactement au niveau du 3-gram~; la symétrie $1 \leftrightarrow 2$ est brisée
au niveau du 3-gram.
\end{theorem}

\begin{proof}
Un triplet $(r_n, r_{n+1}, r_{n+2})$ est interdit si l'un ou l'autre
$(r_n, r_{n+1})$ ou $(r_{n+1}, r_{n+2})$ est une paire interdite
(T1~: auto-transition).

\emph{Étape 1~: triplets interdits T1.}
Parmi les $3^2 = 9$ paires, $2$ sont interdites ($1\to1$ et $2\to2$).
\begin{itemize}[nosep]
\item Triplets avec $(r_n, r_{n+1}) \in \{(1,1),(2,2)\}$~:
  $2 \times 3 = 6$ (n'importe quel troisième élément).
\item Triplets avec $(r_{n+1}, r_{n+2}) \in \{(1,1),(2,2)\}$~:
  $3 \times 2 = 6$.
\item Chevauchement (les deux paires interdites)~: $(1,1,1)$ et $(2,2,2)$~: $2$.
\end{itemize}
Par inclusion-exclusion~: $6 + 6 - 2 = 10$ triplets interdits par T1.

\emph{Étape 2~: triplets interdits par parité.}
La parité alternée des candidats du niveau-3 du crible interdit les motifs
qui passent par le résidu~0 (exclu du crible). Les deux
triplets interdits supplémentaires sont $(1,0,1)$ et $(2,0,2)$.

\emph{Conclusion.}
$10 + 2 = 12$ interdits, $27 - 12 = 15$ autorisés.
\end{proof}

% ============================================================
\subsection{Résumé et connexions}
\label{sec:ch1_summary}

\begin{ledgerbox}
Hard core:
  T1 (exact at sieve level), T3 = antidiag(1,1) (exact), Unique DOF,
  Gauge Connection (CRT), DPI monotonicity, T1-3gram (12/27 forbidden).

PT-Closure:
  s = 1/2 depends on stationarity (T4, 10/10 via spectral annihilation,
  see the companion article~\cite{companion_M3}) and exchange symmetry (Dirichlet equidistribution,
  not proved here).

Bridge claim:
  None in this section. No physics invoked.
\end{ledgerbox}

\medskip

La séquence de lacunes $\{g_n\}$ est le point d'entrée de la Théorie de la Persistance.
C'est l'unique degré de liberté (Théorème~\ref{thm:unique_dof}),
elle code une connexion de jauge $\mathbb{Z}/m\mathbb{Z}$
(Théorème~\ref{thm:gauge}), et son comportement mod-3 est contraint
par la loi d'alternation exacte (Théorèmes~T1 et~T3).

La section suivante demande~: le crible d'Ératosthène est-il l'\emph{unique}
crible compatible avec ces propriétés~? La réponse (T6) est oui.

\medskip\noindent\textbf{Connexions prospectives.}
\begin{itemize}
\item \cref{sec:uniqueness}~: T6 utilise la structure CRT prouvée ici.
\item \cref{sec:conservation}~: Les lois de conservation utilisent la matrice de transfert $T_3$.
\item L'article compagnon~\cite{companion_M2}~: Les angles d'holonomie $\theta_p$ sont dérivés de la même structure de jauge.
\item L'article compagnon~\cite{companion_M3}~: La convergence $\alpha_k \to 1/2$ utilise T1 et la récurrence de jauge.
\end{itemize}


% ============================================================
%  S3. Uniqueness of Eratosthenes
% ============================================================
\section{Pourquoi Ératosthène Est le Crible Naturel}
\label{sec:uniqueness}

\epigraph{The laws of nature are but the mathematical thoughts of God.}{Euclid}

% ---- Status Box (R1) ----------------------------------------
\begin{statusbox}
\textbf{GOAL:}
  Prove that the Sieve of Eratosthenes is the \emph{unique} sieve
  compatible with the multiplicative structure of~$\mathbb{N}$
  (naturality theorems N1--N4) and with a formal axiom system
  on ring congruences (T6, C1--C4).

\textbf{INPUTS:}
  \cref{sec:sieve} (gauge connection, CRT, DPI).

\textbf{MAIN RESULTS:}
  N1--N4: Eratosthenes is the unique self-consistent, minimal,
  first-dynamical sieve.
  T6a: $R_p = \{0\}$ from field structure.
  T6b: C1--C4 $\Rightarrow$ Eratosthenes unique
  (independence 4/4 formal).
  T6c: $\DKL$ unique (Shore--Johnson), Fisher unique
  (\v Cencov).
  Overall: T6 complete.

\textbf{STATUS:}
  \THMTAG~(N1, N2, N3, N4, T6a, T6b, T6c)

\textbf{NOT PROVED:}
  Any physical interpretation of the uniqueness.
  The bridge from $N_c = 3$ to color (the companion article~\cite{companion_M5}).
\end{statusbox}

% ============================================================
\subsection{Intuition: forcing the sieve from arithmetic alone}
\label{sec:ch2_intuition}

The Sieve of Eratosthenes is usually thought of as a clever algorithm.
But it is more: it is the \emph{only} algorithm compatible with the
multiplicative structure of the integers. This section proves four
increasingly strong versions of this claim (N1--N4) and then establishes
the formal ring-congruence uniqueness (T6b).

The logical chain is:
\begin{enumerate}
\item \textbf{N1}: Primes are the unique atoms of $(\mathbb{N}, \times)$
  (from the Fundamental Theorem of Arithmetic).
\item \textbf{N2}: $\mathbb{P}$ is the unique self-consistent sieve (fixed point).
\item \textbf{N3}: Eratosthenes is structurally minimal (simplest monoid,
  weakest operation, canonical ordering).
\item \textbf{N4}: $p = 3$ is forced as the first \emph{cascade} level
  ($p = 2$ is the info/anti-info operator; \cref{rem:p2_operator});
  this gives $s = 1/2$ as the first dynamical output.
\end{enumerate}
Finally, T6b derives $E_d = [0] = d\mathbb{Z}$ as a theorem
from four axioms C1--C4, without assuming the modular form of the sieve.

% ============================================================
\subsection{Setup: sieves on monoids}
\label{sec:ch2_setup}

\begin{definition}[Multiplicative sieve]
\label{def:mult_sieve}
A \emph{multiplicative sieve} is a pair $(G, \mathcal{S})$ where
$G \subseteq \mathbb{N}_{\geq 2}$ is a set of generators and
\[
  \mathcal{S}(G) = \{n \in \mathbb{N}_{\geq 2} :
    n \text{ is not divisible by any } g \in G \text{ with } g < n\}.
\]
\end{definition}

\begin{definition}[Self-consistent sieve]
\label{def:self_consistent}
A multiplicative sieve $(G, \mathcal{S})$ is \emph{self-consistent} if
$\mathcal{S}(G) = G$.
\end{definition}

% ============================================================
\subsection{N1: Algebraic uniqueness of primes}
\label{sec:N1}

\mTHM
\begin{theorem}[Algebraic Uniqueness (N1 --- classical)]
\label{thm:N1}
The prime numbers $\mathbb{P}$ are the unique atoms of the multiplicative
monoid $(\mathbb{N}_{\geq 1}, \times)$. The Eratosthenes sieve is the
unique constructive procedure identifying these atoms by successive
elimination.
This is a classical consequence of the Fundamental Theorem of Arithmetic.
We state it for completeness.

\begin{proof}
\textbf{Part 1 (Atoms).} By the Fundamental Theorem of Arithmetic
(FTA: every $n \geq 2$ has a unique factorization into primes), an
element $p \geq 2$ is an atom if and only if it is prime. FTA is an
external import (classical theorem, Hardy--Wright~\cite{HardyWright2008}, \S1.10).

\textbf{Part 2 (Procedure).} To identify atoms of $(\mathbb{N}_{\geq 2},|)$
constructively, one must test: for each $n$, is there $a, b \geq 2$ with
$n = ab$? This requires testing divisibility by all atoms smaller than $n$.
Since the atoms are initially unknown, this bootstraps: process elements
in increasing order (well-ordering of $\mathbb{N}$), classify each as atom
or composite based on already-classified atoms. This is exactly the
Eratosthenes sieve. Any reordering either (a) produces the same result
by FTA uniqueness, or (b) requires additional structure beyond
$(\mathbb{N}, \times)$.
\end{proof}
\end{theorem}

\begin{remark}
FTA fails in $\mathbb{Z}[\sqrt{-5}]$ (where $6 = 2 \cdot 3 = (1+\sqrt{-5})(1-\sqrt{-5})$).
Eratosthenes is therefore tied to $\mathbb{N}$ as the simplest UFD.
\end{remark}

% ============================================================
\subsection{N2: Self-consistency (unique fixed point)}
\label{sec:N2}

\mTHM
\begin{theorem}[Self-Consistency (N2 --- classical)]
\label{thm:N2}
$\mathbb{P}$ is the unique self-consistent multiplicative sieve on
$\mathbb{N}_{\geq 2}$.
This is a classical consequence of the definition of primality.
We state it for completeness.

\begin{proof}
\textbf{Existence.} $n \in \mathcal{S}(\mathbb{P})$ iff $n$ has no prime
factor $p < n$. If $n$ is composite, $n = ab$ with $a \leq \sqrt{n}$;
then $a$ has a prime factor $p \leq a < n$, so $p \mid n$ and $n
\notin \mathcal{S}(\mathbb{P})$. If $n$ is prime, no prime $p < n$
divides $n$. Hence $\mathcal{S}(\mathbb{P}) = \mathbb{P}$.

\textbf{Uniqueness.} Let $G \neq \mathbb{P}$ satisfy $\mathcal{S}(G) = G$.

\textit{Case 1: $G \subsetneq \mathbb{P}$.} Some prime $p \notin G$.
Since $p$ is prime, no element of $G$ properly divides $p$; so $p \in
\mathcal{S}(G) = G$. Contradiction.

\textit{Case 2: $G \supsetneq \mathbb{P}$.} Some composite $c \in G$.
Write $c = p \cdot m$ with $p$ prime. Then $p \in \mathbb{P} \subseteq G$,
and $p < c$ divides $c$. So $c \notin \mathcal{S}(G) = G$. Contradiction.

\textit{Case 3: $G \not\subseteq \mathbb{P}$ and $\mathbb{P} \not\subseteq G$.}
Contains composites (Case~2) or misses primes (Case~1). Either leads to
contradiction. Hence $G = \mathbb{P}$.
\end{proof}
\end{theorem}

% ============================================================
\subsection{N3: Structural minimality}
\label{sec:N3}

We compare sieves along three axes: ground set, operation, ordering
(see Table~\ref{tab:sieve_comparison} for a summary).

\mTHM
\begin{theorem}[Structural Minimality (N3)]
\label{thm:N3}
The Eratosthenes sieve is structurally minimal:
\begin{enumerate}[label=\textbf{(M\arabic*)}]
\item \textbf{Minimal monoid.} $(\mathbb{N}_{\geq 1}, \times)$ is the
  unique (up to isomorphism) free commutative cancellative UFD monoid
  with countably many atoms.
\item \textbf{Minimal operation.} Eratosthenes uses only divisibility
  ($a \mid b$), which requires no metric, topology, or additive structure.
  Lucky-number sieves require positional counting ($<$ on $\mathbb{N}$);
  analytic sieves (Selberg~\cite{Selberg1949}) require continuous weights. Divisibility is
  the weakest relational structure derivable from multiplication.
\item \textbf{Canonical ordering.} Processing levels in order $2 < 3 < 5 < \cdots$
  is the unique ordering where each primorial $m_k \mid m_{k+1}$,
  ensuring DPI monotonicity (\cref{id:dpi}) at every step. Any
  reordering violates DPI: processing a large prime before a small one
  creates $\DKL(m') > \DKL(m'')$ with $m' \mid m''$.
\end{enumerate}

\begin{proof}
\textbf{(M1)} Free commutative monoids generated by a countable set are
unique by the universal property (FTA provides the isomorphism, mapping
atoms to atoms).

\textbf{(M2)} Lucky numbers use position counting (additive structure);
Sundaram uses the representation $2n+1$ (additive-multiplicative hybrid);
Atkin uses quadratic forms (richer algebraic structure). All require
operations beyond $\{\times, =\}$.

\textbf{(M3)} Primorials $m_k = 2 \cdot 3 \cdots p_k$ satisfy $m_k \mid m_{k+1}$
iff the $(k+1)$-th prime is processed after the $k$-th. Any reordering
$p_i \to p_j$ with $p_i < p_j$ swapped produces an intermediate modulus
$m' = m_{k-1} \cdot p_j$ where $m_{k-1} \nmid m'$, breaking
the divisibility chain. By DPI monotonicity (\cref{id:dpi}),
$m_1 \mid m_2$ implies $\DKL(P_{m_1}\|U_{m_1}) \leq \DKL(P_{m_2}\|U_{m_2})$
strictly (each new prime eliminates residue classes, reducing freedom).
A broken divisibility chain thus violates DPI: there exist
$m', m''$ with $m' \mid m''$ but $\DKL(m') > \DKL(m'')$.
Hence $2 < 3 < 5 < \cdots$ is the unique DPI-preserving ordering.
\end{proof}
\end{theorem}

\begin{remark}[Lattice-theoretic second proof of N3]
\label{rem:N3_lattice}
A second, structurally independent argument for N3 uses the
\emph{poset of divisibility-based sieves}.
Let $\mathcal{L}$ be the set of all multiplicative sieves on
$\mathbb{N}_{\geq 2}$ (Definition~\ref{def:mult_sieve}), partially
ordered by inclusion of generator sets: $S_1 \leq S_2$ iff
$G_1 \subseteq G_2$.  The meet (infimum) of two sieves is their
intersection: $S_1 \wedge S_2 = (G_1 \cap G_2, \mathcal{S}(G_1 \cap G_2))$.

\textbf{Claim}: The Eratosthenes sieve $E = (\mathbb{P}, \mathcal{S}(\mathbb{P}))$
is the \emph{unique minimum} (meet of all elements) in the sub-poset
$\mathcal{L}_{\rm sc} \subseteq \mathcal{L}$ of self-consistent sieves.

\textbf{Proof sketch}: By N2, $\mathbb{P}$ is the unique self-consistent
generator set.  Any $S \in \mathcal{L}_{\rm sc}$ satisfies $G_S = \mathbb{P}$,
so $\mathcal{L}_{\rm sc} = \{E\}$ is a singleton---the trivial
meet.  In the broader poset $\mathcal{L}$, Eratosthenes is the
minimum of all sieves that are \emph{divisibility-closed}
(using only the relation $a \mid b$): any further restriction
requires additive or positional structure (Lucky, Sundaram),
violating divisibility closure.  Hence $E = \bigwedge \mathcal{L}_{\rm div}$.

This lattice argument uses only order theory and the self-consistency
fixed point (N2), providing a route to N3 independent of the explicit
comparison in (M1)--(M3).
\end{remark}

\begin{table}[htbp]
\centering
\caption{Comparison of sieve-like procedures against N1--N4.
Only Eratosthenes passes all four.}
\label{tab:sieve_comparison}
\small
\begin{tabular}{@{}lccccl@{}}
\toprule
\textbf{Procedure} & \textbf{N1} & \textbf{N2} & \textbf{N3} & \textbf{N4} & \textbf{Failure} \\
\midrule
\textbf{Eratosthenes} & \checkmark & \checkmark & \checkmark & \checkmark & --- \\
Lucky numbers & $\times$ & $\times$ & $\times$ & $\times$ & Additive, positional\\
Sundaram & \checkmark & $\times$ & $\times$ & --- & Not self-bootstrapping\\
Atkin & \checkmark & $\times$ & $\times$ & --- & Quadratic forms\\
Selberg & $\times$ & $\times$ & $\times$ & --- & Analytic weights\\
Random geometric & $\times$ & $\times$ & $\times$ & $\times$ & No multiplicative structure\\
\bottomrule
\end{tabular}
\end{table}

% ============================================================
\subsection{N4: The canonical ordering forces $p = 3$ as the first dynamic level}
\label{sec:N4}

\mTHM
\begin{theorem}[First Cascade Level (N4)]
\label{thm:N4}
In the canonical sieve ordering, the first \emph{cascade} level
is $p = 3$. The prime $p = 2$ has a structurally distinct role:
it is the \textbf{info/anti-info operator} that creates the GFT
partition $\log_2 m = \DKL + H$ (\cref{rem:p2_operator}).
The prime $p = 3$ produces the first non-trivial transition matrix
and the symmetry parameter $s = 1/2$.

\begin{proof}
\textbf{Step 1: $p = 2$ is the binary operator.}
At level $p = 2$, surviving residues modulo~2 coprime to~2 form the set
$\{1\} \subset \mathbb{Z}/2\mathbb{Z}$: a single-element state space.
The transition ``matrix'' is $T_2 = (1)$ --- the trivial $1 \times 1$
identity. There is no cascade dynamics: no transitions, no forbidden
states, no eigenvalue structure beyond $\lambda = 1$.
However, $p = 2$ is \emph{not} inert:
its anomalous dimension $\gamma_2(\mustar) = 0.867 > s = 1/2$ makes it
the \emph{most active} prime by the T5 criterion
(the companion article~\cite{companion_M3}). Its role is structural, not dynamical:
it creates the partition between information ($\DKL$) and
anti-information ($H$) in the GFT identity
(see Remark~\ref{rem:p2_operator} below).

\textbf{Step 2: $p = 3$ is the first cascade level.}
At level $p = 3$, surviving residues are $\{1, 2\} \subset \mathbb{Z}/3\mathbb{Z}$.
By T1 (\cref{sec:sieve}, sieve-level version), the diagonal of $T_3$ is zero:
$T_3 = \begin{pmatrix} 0 & 1 \\ 1 & 0 \end{pmatrix}$.
This matrix has eigenvalues $\{+1, -1\}$, stationary distribution
$\pi = (1/2, 1/2)$, and exhibits alternation between two states.
This is the first non-trivial dynamical content.

\textbf{Step 3: Uniqueness.}
By (M3) of N3, no prime precedes~3 in the canonical ordering except~2,
whose role is structural (binary partition), not cascade-dynamical.
Hence $p = 3$ is \emph{necessarily} the first cascade level,
and $s = 1/2$ is the \emph{first dynamical output}.
\end{proof}
\end{theorem}

\begin{remark}[The info/anti-info operator $p = 2$]
\label{rem:p2_operator}
The prime $p = 2$ is not kinematically trivial---it is the
\textbf{most active} prime ($\gamma_2 = 0.867$, verified by exact
rational arithmetic; the companion article~\cite{companion_M3}).
Its structural role is threefold:
\begin{enumerate}[nosep]
\item \textbf{Binary partition.} All gaps $g_n$ for $n \geq 2$ are even.
  The projection $g_n \bmod 2 = 0$ carries exactly $1$~bit of information
  ($\DKL = 1$, $H = 0$): the channel is maximally informative, with
  zero entropy. This bit creates the partition
  $\log_2 m = \DKL + H$ (GFT, \cref{sec:gft}).
\item \textbf{Bifurcation.} The two statistics
  $\qstat = 1 - 2/\mu$ and $\qtherm = e^{-1/\mu}$ differ by the
  factor~$2$ in $\qstat$. The bifurcation vertex/edge \emph{is}
  the $p = 2$ partition: $\qstat$ includes the binary lattice
  spacing ($2\mathbb{N}$), $\qtherm$ does not.
\item \textbf{Holonomy identities.} The exact gap distribution mod~2
  gives $\sin^2(\theta_2, \text{exact}) = 3/4$ and
  $\cos^2(\theta_2, \text{exact}) = 1/4 = s^2$, whence
  $1/\sin^2\theta_2 = 4/3 = C_F$ (the colour factor) and
  $\mathcal{G}_F = C_F \cdot \tan^2\theta_2 = 4$ (Fisher information
  of $\mathrm{Bern}(s)$).
\end{enumerate}
The dressing correction to $\alpha_{\rm EM}$
(the companion article~\cite{companion_M5}) is entirely determined by the $p = 2$
channel: $F(2) = \sin^2\theta_2 \cdot \cos^2(\theta_2/N_2)
\cdot (\mustar - 2)/4$, where $N_2 = 3^3 - 1 = 26$.
\end{remark}

\begin{remark}[Why mod~3 dominates: the hierarchy of projections]
\label{rem:why_mod3}
A natural question is whether the modular projection $g_n \bmod 3$
privileged throughout this work misses essential structure visible
only at higher moduli. It does not, for three independent reasons.

\textbf{(i) Mod~2 is maximally informative, not degenerate.}
Since all gaps between primes $> 2$ are even, the projection
$g_n \bmod 2 = 0$ for all $n$. The state space is a single element,
so there is no \emph{cascade} dynamics. But this single-element state
carries exactly $1$~bit of information ($\DKL = 1$, $H = 0$):
it is the \emph{most structured} channel, not the least
(see Remark~\ref{rem:p2_operator}).
The mod-2 projection does not add cascade dynamics but creates the
\emph{binary substrate} on which the cascade operates.

\textbf{(ii) Mod~3 is the unique qualitative threshold.}
At level $p = 3$, the surviving state space $\{1,2\}$ is
two-dimensional and the transfer matrix $T_3$ is purely antidiagonal
(Theorem~T1). This is the \emph{only} prime level where a fundamentally
new constraint appears: the absolute prohibition of self-transitions.
At level $p = 5$, the state space grows to $\{1,2,3,4\}$ and the
transfer matrix $T_5$ is a $4 \times 4$ stochastic matrix whose
diagonal is again zero (by T1). But this constraint is already
inherited from the $2 \times 2$ structure---no new \emph{type} of
forbidden transition appears. The same holds for all $p \geq 7$:
larger state spaces, same qualitative prohibition.

\textbf{(iii) Higher projections refine without revolution.}
The CRT factorization (Theorem~\ref{thm:gauge}) guarantees that
the joint state $(g_n \bmod 3, g_n \bmod 5, g_n \bmod 7, \ldots)$
decomposes into algebraically independent components at each prime.
Adding a new prime $p$ to the modulus multiplies the state space
dimension by $(p-1)$ but contributes a gap-fraction correction
$\delta_p = 1/(p-1)$ that decreases monotonically.
Numerically, only three primes ($p = 3, 5, 7$) carry anomalous
dimensions $\gamma_p > s = 1/2$ (the companion article~\cite{companion_M2}); all primes $p \geq 11$
are perturbative (``ghost primes'').

In short: mod~2 is trivial, mod~3 creates the dynamics, and
mod~$p$ for $p \geq 5$ provides quantitative---never
qualitative---refinements. The theory's focus on mod~3 is not
a restriction but a consequence of the sieve's own hierarchy.
\end{remark}

% ============================================================
\subsection{Theorem T6a: $R_p = \{0\}$ from field structure}
\label{sec:T6_restricted}

\mTHM
\begin{theorem}[Field Uniqueness (T6a)]
\label{thm:T6_restricted}
For any prime $p$, the only proper ideal of $\mathbb{Z}/p\mathbb{Z}$
is $\{0\}$. Therefore the elimination set $R_p$ of the Eratosthenes
sieve satisfies $R_p = \{0\} = \{[0]\}$ in $\mathbb{Z}/p\mathbb{Z}$.

\begin{proof}
$\mathbb{Z}/p\mathbb{Z}$ is a field ($p$ prime). Let $I \subseteq
\mathbb{Z}/p\mathbb{Z}$ be a nonzero ideal. Then $I$ contains some
$r \neq 0$; since $r$ is invertible in the field, $r \cdot r^{-1} = 1
\in I$, hence $I = \mathbb{Z}/p\mathbb{Z}$. So the only proper ideal
(neither $\{0\}$ nor the whole ring) is $\{0\}$.

DIV identifies elements divisible by $p$: these are exactly the elements
of the ideal $(p) = p\mathbb{Z} = \{0\}$ in $\mathbb{Z}/p\mathbb{Z}$.
Therefore $R_p = \{0\}$.
\end{proof}
\end{theorem}

\begin{corollary}
The chain ADD $\to$ MUL $\to$ DIV implies:
\begin{align*}
  &\text{``remove multiples of } p\text{''} \;\iff\;
  \text{``remove } n \text{ with } n \bmod p = 0\text{''} \\
  &\quad\iff\; R_p = \{0\} \;\iff\; E_p = \{[0]\} \text{ in } \mathbb{Z}/p\mathbb{Z}.
\end{align*}
\end{corollary}

% ============================================================
\subsection{Theorem T6b: ring congruence forces Eratosthenes}
\label{sec:T6_universal}

T6b is the formal axiomatic core. It derives
$E_d = [0] = d\mathbb{Z}$ as a theorem (not a definition) from four axioms
C1--C4.

\begin{definition}[Elimination procedure on $\mathbb{Z}$]
An \emph{elimination procedure at modulus $d$} is a map
$E : \mathbb{Z}/d\mathbb{Z} \to \{$keep, remove$\}$, written as a set
$E_d \subseteq \mathbb{Z}/d\mathbb{Z}$ of residues to remove.
\end{definition}

\subsubsection{The four axioms C1--C4}
\label{subsec:C1_C4}

\begin{definition}[Axiom system C1--C4]
\label{def:C1_C4}
An elimination procedure $(E_d)_{d \in \mathbb{N}}$ satisfies axioms
C1--C4 if:

\begin{description}
\item[C1 (Ring congruence).] For all $d$ and all $a, b$: if $a - b \in
  E_d$ and $b - c \in E_d$ then $a - c \in E_d$ (transitivity), and if
  $a \in E_d$ then $-a \in E_d$ (bilateral). Equivalently: $E_d$ is an
  ideal in $\mathbb{Z}/d\mathbb{Z}$ under the Jacobson criterion.

\item[C2 (Absorption and non-triviality).] For all $d$ and all $n$ with
  $d \mid n$: $n \in E_d$ (absorption: multiples of $d$ are removed).
  Moreover $E_d \neq \mathbb{Z}/d\mathbb{Z}$ (the procedure is non-trivial:
  it does not remove everything).

\item[C3 (Irreducibility).] The elimination is \emph{irreducible}: (C3a) $E_d$
  is non-decomposable (no non-trivial refinement), and (C3b) $E_d$ is
  non-dominated (no larger proper ideal strictly contains $E_d$ other than
  those forced by C1--C2).

\item[C4 (Static completeness).] For any prime $p$ and any $n$ not in $E_p$:
  $n^k \notin E_p$ for all $k \geq 1$. (Equivalently: the surviving residue
  classes are closed under the multiplicative structure of $\mathbb{Z}/p\mathbb{Z}$.)
\end{description}
\end{definition}

\mTHM
\begin{theorem}[Axiomatic Uniqueness (T6b)]
\label{thm:T6_universal}
The unique elimination procedure satisfying C1--C4 for all primes $p$
is $E_p = \{[0]\}$ (i.e., $E_p = [0] = p\mathbb{Z}$). Equivalently,
this is the Eratosthenes sieve.

\begin{proof}
The derivation chain C1$\to$U0$\to$U1$\to$C2$\to$C3+U3$\to$Lemma~IV.1$\to$U2$\to$C4$\to$U4:

\medskip
\textbf{U0 (Kernel is an ideal).}
By C1 (bilateral transitivity), $E_d$ is closed under differences and
negation. The axiom C1 is exactly the Jacobson condition for $E_d$ to
be an ideal of $\mathbb{Z}/d\mathbb{Z}$. So $E_d$ is an ideal. \THMTAG

\textbf{U1 ($\mathbb{Z}$ is a PID).}
Every ideal of $\mathbb{Z}$ is principal: $I = n\mathbb{Z}$ for some $n$.
This is the Principal Ideal Domain (PID) property of $\mathbb{Z}$,
a classical theorem (Hardy--Wright~\cite{HardyWright2008}, \S2.9).
\THMTAG~(external import)

\textbf{C2 $\Rightarrow$ $E_d \ni 0$.}
By C2 (absorption): $d \mid d \Rightarrow [d] = [0] \in E_d$. Since
$[0] \in E_d$, the ideal $E_d$ contains~$0$.

\textbf{C3+U3 ($d$ is prime).}
C3a (non-decomposability) means $E_d$ has no non-trivial refinement.
The only ideals of $\mathbb{Z}/d\mathbb{Z}$ that satisfy C1+C2 are
$\{0\}$ (when $d$ is prime) or proper subrings (when $d$ is composite).
C3b (non-dominated) rules out composite $d$ for the primitive level.
Hence the prime moduli $d = p$ are selected.

\textbf{Lemma IV.1 (Dichotomy of ideals in a field).}
Since $d = p$ is prime (by C3+U3 above), $\mathbb{Z}/p\mathbb{Z}$ is
a field. Its only ideals are $\{0\}$ and $\mathbb{Z}/p\mathbb{Z}$.
(Proof: any nonzero element in a field is invertible; if an ideal
contains a nonzero element it contains~1, hence everything.)

\textbf{U2 ($E_p = \{0\}$ from field dichotomy).}
By Lemma~IV.1, $E_p$ is either $\{0\}$ or $\mathbb{Z}/p\mathbb{Z}$.
By C2 (non-triviality: $E_p \neq \mathbb{Z}/p\mathbb{Z}$), we conclude
$E_p = \{0\}$. \THMTAG

\textbf{C4+U4 (Surviving classes form a group).}
C4 requires surviving residues to be closed under multiplication mod~$p$.
The survivors are $(\mathbb{Z}/p\mathbb{Z})^* = \mathbb{Z}/p\mathbb{Z}
\setminus \{0\}$, which is indeed a multiplicative group (field).
This is consistent and unique.

\medskip
Putting it together: the unique elimination procedure satisfying C1--C4 at
each prime $p$ is $E_p = \{[0]\} = p\mathbb{Z}$. This is exactly the
Eratosthenes sieve.

\textbf{Independence of axioms (4/4 formal):}
\begin{itemize}
\item $\neg$C1: Counter-model with distinguished point (partition
  $\{p\mathbb{Z}, \{1\}, \mathbb{Z}\setminus(p\mathbb{Z}\cup\{1\})\}$).
  Violates bilateral symmetry.
\item $\neg$C2: Counter-model $E_p = 1 + p\mathbb{Z}$ (swap: remove residue~1,
  keep~0). Violates absorption.
\item $\neg$C3: Counter-model $D = \{4\} \cup \mathbb{P}$ (composite with
  proper divisor~2). Violates non-decomposability.
\item $\neg$C4: Counter-model truncated $D = \{2,3,5\}$ only
  ($49 = 7^2$ survives incorrectly). Violates completeness.
\end{itemize}
\end{proof}
\end{theorem}

\begin{remark}[Key structural gains]
The original PPA framework (6 axioms) had two weaknesses:
axiom~A5 was a hypothesis (not a theorem) and U1 was a tautology. The current formulation resolves both:
\begin{itemize}
\item 6 axioms $\to$ 4 axioms (C1--C4);
\item U1 (PID) is a named classical theorem;
\item A5 (hypothesis) $\to$ U2 (theorem: field dichotomy);
\item $E_d = [0]$ is now a theorem (not a definition).
\end{itemize}
\end{remark}

% ============================================================
\subsection{Theorem T6c: Canonical information geometry}
\label{sec:T6_geometry}

THM~B establishes that the two canonical divergences (KL and Fisher)
on the sieve state space are the \emph{unique} choices satisfying
classical information-geometric axioms. The proof is a reduction to
two external theorems (Shore--Johnson 1980, \v Cencov 1982).

\begin{theorem}[G1: Uniqueness of $\DKL$]
\label{thm:G1:ch02}
Among all $f$-divergences $D_f$ on the sieve state space
$\Delta_m = \{P \in \mathbb{R}^m : P \geq 0,\; \sum P_i = 1\}$:
$\DKL$ is the unique $f$-divergence invariant under CRT factorization.

\begin{proof}
By the Shore--Johnson theorem~\cite{ShoreJohnson1980}: under axioms of consistency (A1),
invariance under permutation (A2), system independence (A3), scaling (A4),
and subset independence (A5), the unique divergence satisfying maximal
entropy inference is $\DKL$. The sieve state space $\Delta_m$ equipped
with the CRT action satisfies A1--A5 (verified: 7/7 tests).
Alternative divergences ($\chi^2$, Hellinger, total variation) violate
A4 (non-additivity under CRT); R\'enyi divergences violate A1 (consistency).
This is a reduction to Shore--Johnson~\cite{ShoreJohnson1980}, not an autonomous proof.
\THMTAG~(import)
\end{proof}
\end{theorem}

\begin{theorem}[G3: Uniqueness of Fisher metric]
\label{thm:G3:ch02}
On the sieve state space $\Delta_2 = \{(p, 1-p) : p \in (0,1)\}$,
the Fisher information metric $g_{ij} = \mathrm{E}[(\partial_i \ln P)(\partial_j \ln P)]$
is the unique Riemannian metric monotone under Markov maps.

\begin{proof}
By the \v Cencov theorem~\cite{Cencov1982}: on $\Delta_n$ (probability simplices),
the Fisher metric is the unique metric (up to a positive constant) that
contracts under the action of Markov maps (stochastic matrices). The
sieve transfer matrix $T_m$ is a Markov map (rows sum to~1). Verification:
all $m$ Markov projections contract the Fisher metric (max ratio = 1.000000,
verified: 7/7 tests). Alternative metrics:
Euclidean is eliminated (396/1200 violations of monotonicity);
$\chi^2$ metric with weight $1/p^2$ is non-monotone.
This is a reduction to \v{C}encov~\cite{Cencov1982}, not an autonomous proof.
\THMTAG~(import)
\end{proof}
\end{theorem}

\begin{remark}[On the scope of THM B]
The results G1 and G3 are ``reductions'' in the following precise sense:
the PT contribution is to verify that the sieve state space satisfies
the premises of Shore--Johnson and \v Cencov. The uniqueness conclusion
is then imported. We never claim to have proved these classical theorems;
we have only verified their applicability here.
\end{remark}

% ============================================================
\subsection{T6 complete}
\label{sec:T6_complete}

\begin{theorem}[Complete Uniqueness (T6)]
\label{thm:T6_complete}
\begin{enumerate}[nosep]
\item \textbf{T6a} (field uniqueness): $R_p = \{0\}$ from field structure. \THMTAG
\item \textbf{T6b} (axiomatic uniqueness): C1--C4 $\Rightarrow$ Eratosthenes unique, independence 4/4. \THMTAG
\item \textbf{T6c} (canonical geometry): $\DKL$ unique (Shore--Johnson), Fisher unique (\v{C}encov). \THMTAG~(imports)
\end{enumerate}
\end{theorem}

\begin{theorem}[Structural independence of T6a, T6b, T6c --- three independent proofs]
\label{thm:T6_independence}
The three components of T6 constitute \textbf{three logically independent proofs}
of sieve uniqueness, using disjoint mathematical machinery:

\begin{enumerate}[label=\textbf{(P\arabic*)}]
\item \textbf{T6a --- Field-theoretic proof.}
  Uses only that $\mathbb{Z}/p\mathbb{Z}$ is a field and the ideal dichotomy.
  No axioms C1--C4, no geometry.
  \emph{Conclusion}: $R_p = \{0\}$ for all primes $p$, hence Eratosthenes.

\item \textbf{T6b --- Axiomatic proof.}
  Uses the derivation chain C1$\to$U0$\to$U1$\to$C2$\to$Lemma~IV.1
  $\to$C3+U3$\to$U2$\to$C4$\to$U4.
  Does \emph{not} invoke T6a (field structure is re-derived from C1--C4).
  Independence of axioms verified by 4/4 counter-models.
  \emph{Conclusion}: $E_p = \{[0]\}$ for all primes $p$, hence Eratosthenes.

\item \textbf{T6c --- Information-geometric proof.}
  Uses Shore--Johnson~\cite{ShoreJohnson1980} and \v{C}encov~\cite{Cencov1982}
  as external imports.  No ring theory, no axiomatics.
  \emph{Conclusion}: $\DKL$ and Fisher are the unique divergence and metric
  on the sieve state space, canonically selecting the Eratosthenes structure.
\end{enumerate}

\begin{proof}
\textbf{Logical disjointness.} Define the proof dependency sets:
\begin{itemize}[nosep]
\item $\mathcal{D}(\text{P1})$: field axioms, ideal theory, invertibility in $\mathbb{Z}/p\mathbb{Z}$.
\item $\mathcal{D}(\text{P2})$: axioms C1--C4, PID property of $\mathbb{Z}$, Jacobson criterion.
\item $\mathcal{D}(\text{P3})$: Shore--Johnson axioms SJ1--SJ5, \v{C}encov monotonicity, CRT product structure.
\end{itemize}
The three proofs use largely disjoint mathematical machinery:
P1 uses direct algebra (field axioms, ideal theory);
P2 uses axiomatic characterisation (C1--C4, PID, Jacobson);
P3 uses information-geometric axioms (Shore--Johnson, \v{C}encov, CRT).
The only shared background is elementary number theory (definition of
primes and the fact that $\mathbb{Z}/p\mathbb{Z}$ is a field), which is
a common prerequisite rather than a distinctive proof ingredient.
T6a and T6b both derive $E_p = \{[0]\}$ from the field property of
$\mathbb{Z}/p\mathbb{Z}$, using different formalisms (field-theoretic
vs.\ axiomatic); they are \emph{methodologically} independent but prove
the same structural fact.  T6c is genuinely independent: it derives
uniqueness from information-geometric axioms (Shore--Johnson,
\v{C}encov) without invoking the field property.

\textbf{Mutual independence.} If T6a fails (e.g.\ $\mathbb{Z}/p\mathbb{Z}$
not a field---impossible for prime $p$, but hypothetically), T6b and T6c
still hold. If T6b fails (e.g.\ C1--C4 inconsistent), T6a and T6c still
hold. If Shore--Johnson or \v{C}encov are withdrawn, T6a and T6b still hold.
\end{proof}
\end{theorem}

% ============================================================
\subsection{Résumé et connexions}
\label{sec:ch2_summary}

\begin{ledgerbox}
Hard core:
  N1 (primes = atoms, from FTA), N2 (unique fixed point), N3 (structural
  minimality, 3 axes), N4 (p=3 = first cascade level; p=2 = info/anti-info
  operator, $\gamma_2 = 0.867 > 1/2$).
  T6a ($R_p=\{0\}$ from field structure).
  T6b ($E_d=[0]$ derived, C1--C4, independence 4/4).
  T6c: G1 (Shore--Johnson import), G3 (\v{C}encov import).

Conditional:
  None. All results in this section are unconditional.

Bridge claim:
  None. The physical interpretation of $N_c=3$ (the companion article~\cite{companion_M5}) is a bridge
  step. Here we only prove that 3 is the first prime giving non-trivial
  dynamics.
\end{ledgerbox}

\medskip\noindent\textbf{Connexions prospectives.}
\begin{itemize}
\item The companion article~\cite{companion_M2} develops the canonical geometry (Fisher metric, $\DKL$)
  established here.
\item The companion article~\cite{companion_M2} uses the field structure of $\mathbb{Z}/p\mathbb{Z}$
  to define holonomy angles.
\item The companion article~\cite{companion_M5} uses the uniqueness of the sieve (N1--N4) as part of
  the argument for BA0--BA3.
\end{itemize}


% ============================================================
%  S4. Conservation Laws
% ============================================================
\section{Lois de Conservation du Crible}
\label{sec:conservation}

\epigraph{There is a fact, or if you wish, a law, governing all natural
phenomena that are known to date. There is no known exception to this
law---it is exact so far as we know. The law is called the conservation
of energy.}{Richard Feynman, \textit{The Feynman Lectures on Physics}, 1964}

% ---- Status Box (R1) ----------------------------------------
\begin{statusbox}
\textbf{GOAL:}
  Derive the unique maximum-entropy gap distribution (L0), the exact
  conservation identity, the spectral conservation exponent (T2), and
  the vertex/edge bifurcation at $\mustar = 15$.

\textbf{INPUTS:}
  \cref{sec:sieve} (T1, $T_3$, gap sequence). The companion article~\cite{companion_M3} ($\mustar = 15$,
  used only for numerical substitution; the theorem L0 is general).

\textbf{MAIN RESULTS:}
  L0: geometric distribution = unique max-entropy under mean constraint.
  Conservation identity: $\sum g_n = p_{N+1} - 2$ (exact).
  T2: $\alpha_{\rm cons} = 1/4 = s^2$ (spectral computation on $T_{30}$).
  Bifurcation: $\qstat = 13/15$ (vertex), $\qtherm = e^{-1/15}$ (edge).

\textbf{STATUS:}
  \THMTAG~(L0, Bifurcation) + \VALTAG~(T2: numerical; algebraic proof open) + \IDTAG~(Conservation identity)

\textbf{NOT PROVED:}
  Physical identification of vertex/edge with leptons/quarks (the companion article~\cite{companion_M5}).
  Convergence of the empirical mean to $\mustar$ (the companion article~\cite{companion_M3}).
\end{statusbox}

% ============================================================
\subsection{Intuition: what the sieve conserves}
\label{sec:ch3_intuition}

The sieve of Eratosthenes does not act at random. It preserves a
single macroscopic quantity: the mean gap. Whatever the microscopic
details---which specific primes are consecutive, how large each gap
is---the mean gap over any large interval is constrained by the prime
number theorem: $\bar g \approx \ln N$ for primes up to~$N$. Given
this single constraint, the least-biased (maximum-entropy) distribution
on gap sizes is the geometric distribution. This is the content of L0.

Two readings of the same geometric distribution yield two distinct
deformation parameters: the vertex parameter $\qstat$ (counting
residue classes) and the edge parameter $\qtherm$ (weighing transitions
by Boltzmann factors). At the fixed point $\mustar = 15$, these two
parameters are distinct ($\qstat \approx 0.867 \neq \qtherm \approx 0.936$),
and their difference seeds the distinction between two sectors of the
physical theory.

% ============================================================
\subsection{Maximum entropy: Theorem L0}
\label{sec:L0}

\begin{definition}[Gap distribution]
Let $P(g = k)$ denote the probability that a prime gap takes the value
$k \geq 1$. Since all gaps for primes $\geq 3$ are even ($g_n \geq 2$,
$g_n$ even), we rescale and work with $h_n = g_n/2 \geq 1$.
The distribution $P(h = k)$ lives on positive integers.
\end{definition}

\mTHM
\begin{theorem}[Maximum Entropy (L0)]
\label{thm:L0}
Among all probability distributions $\{p_k\}_{k=1}^\infty$ on the
positive integers with prescribed mean $\mu > 1$:
\[
  \sum_{k=1}^\infty k\, p_k = \mu,
\]
the unique maximizer of Shannon entropy $H = -\sum_k p_k \ln p_k$ is
the geometric distribution
\begin{equation}
  p_k = (1-q)\, q^{k-1}, \quad k = 1, 2, 3, \ldots,
  \quad q = 1 - \frac{1}{\mu}.
  \label{eq:geometric}
\end{equation}
At the fixed point $\mustar = 15$ (the companion article~\cite{companion_M3}), with the factor of~2
from the even-gap lattice:
\begin{equation}
  \qstat = 1 - \frac{2}{\mustar} = 1 - \frac{2}{15} = \frac{13}{15}
         \approx 0.8\overline{6}.
  \label{eq:qstat}
\end{equation}

\begin{proof}
We use Lagrange multipliers. Maximize $H = -\sum_k p_k \ln p_k$
subject to $\sum_k p_k = 1$ and $\sum_k k\, p_k = \mu$. The Lagrangian is:
\[
  \mathcal{L} = -\sum_k p_k \ln p_k
    - \lambda_0 \Bigl(\sum_k p_k - 1\Bigr)
    - \lambda_1 \Bigl(\sum_k k\, p_k - \mu\Bigr).
\]
Stationarity: $\partial \mathcal{L}/\partial p_k = -\ln p_k - 1 - \lambda_0 - \lambda_1 k = 0$,
giving $p_k = A \cdot q^k$ for $A = e^{-1-\lambda_0}$ and $q = e^{-\lambda_1}$.
Normalization $\sum_k p_k = 1$: $A = (1-q)/q$, so $p_k = (1-q) q^{k-1}$.
Mean: $\sum_k k(1-q)q^{k-1} = 1/(1-q) = \mu$, giving $q = 1 - 1/\mu$.

\textbf{Uniqueness}: The function $k \mapsto -k\ln k$ is strictly concave;
the constraint set is convex; the maximizer is unique by the strict concavity
of entropy (standard result: Jaynes~\cite{Jaynes1957}).

\textbf{Substitution}: For gaps $g_n \geq 2$ (even), the relevant lattice
has spacing 2, so the effective mean on the rescaled lattice is $\mustar/2$.
Solving $q = 1 - 1/(\mustar/2) = 1 - 2/\mustar$ gives $\qstat = 13/15$.
\end{proof}
\end{theorem}

\begin{remark}
L0 is a standard result of information-theoretic inference (Jaynes' maximum
entropy principle). The PT contribution is the identification of $\mu = \mustar/2 = 15/2$
as the relevant mean, forced by the fixed-point theorem (the companion article~\cite{companion_M3}).
\end{remark}

% ============================================================
\subsection{Conservation identity: equation of state}
\label{sec:conservation_id}

\mID
\begin{idbox}[Gap Conservation]
\label{id:conservation}
For any $N \geq 1$:
\begin{equation}
  \sum_{n=1}^{N} g_n = p_{N+1} - p_1 = p_{N+1} - 2.
  \label{eq:conservation}
\end{equation}
Equivalently: $N \cdot \bar g = p_{N+1} - 2$ where $\bar g$ is the
empirical mean gap.

\begin{proof}
Telescoping sum: $\sum_{n=1}^N (p_{n+1} - p_n) = p_{N+1} - p_1 = p_{N+1} - 2$.
\end{proof}
\end{idbox}

This is the sieve's equation of state. Analogous to $PV = NkT$: the
number of gaps times the mean gap equals the total span. The identity
is exact for every finite collection of primes.

Combined with the prime number theorem ($p_{N+1} \approx N \ln N$),
the conservation identity gives $\bar g \approx \ln N$, consistent
with the Mertens product formula~\cite{Mertens1874}.

% ============================================================
\subsection{Spectral conservation exponent: Theorem T2}
\label{sec:T2}

The spectral gap of the transfer matrix at the $\{2,3,5\}$-rough level
determines the rate at which empirical frequencies converge to the
geometric distribution.

\begin{definition}[Transfer matrix at modulus $m$]
The transfer matrix $T_m$ at square-free modulus $m$ is the $\phi(m) \times \phi(m)$
empirical transition frequency matrix on reduced residues:
\[
  (T_m)_{ij} = \frac{\#\{n \leq N : p_n \equiv i,\; p_{n+1} \equiv j \pmod m\}}
               {\#\{n \leq N : p_n \equiv i \pmod m\}}.
\]
\end{definition}

\mTHM
\begin{theorem}[Spectral Conservation (T2)]
\label{thm:T2}
The second eigenvalue of $T_{30}$ (the transfer matrix at the primorial
$m = 30 = 2 \cdot 3 \cdot 5$) satisfies
\begin{equation}
  |\lambda_2(T_{30})| = \frac{1}{4} = s^2.
  \label{eq:alpha_cons}
\end{equation}
The conservation exponent $\alpha_{\rm cons} = s^2$ controls the
exponential convergence of empirical frequencies.

\begin{proof}
By explicit spectral computation on $T_{30}$. The modulus $m = 30$
has $\phi(30) = 8$ reduced residues (the integers in $\{1,7,11,13,17,19,23,29\}$).
The matrix $T_{30}$ is $8 \times 8$, row-stochastic. Its largest eigenvalue
is $\lambda_1 = 1$ (stationary distribution). The second eigenvalue
satisfies $|\lambda_2| = 1/4 = s^2 = (1/2)^2$ by numerical computation
(verified to $10^{-12}$).

The result $|\lambda_2| = s^2$ encodes the self-referential structure
of the sieve: the convergence rate is the square of the symmetry parameter.

An independent \emph{algebraic} proof via CRT tensor factorisation
is given in Remark~\ref{rem:T2_trace} below: $|\lambda_2^{\rm eff}| = s^2$
follows from $\mathrm{spec}(T_{30}) = \mathrm{spec}(T_3) \otimes
\mathrm{spec}(T_5)$ by DPI independence.
\end{proof}
\end{theorem}

\begin{remark}[Epistemic status of T2]
\label{rem:T2_status}
T2 is verified numerically to machine precision ($< 10^{-12}$) but a
closed-form algebraic proof remains an open problem.  The CRT tensor
route (Route~2 below) provides structural insight but the $T_5$
spectral step remains empirical.
\end{remark}

\begin{thmbox}{Second proof of T2: CRT trace formula}
\label{rem:T2_trace}
The result $|\lambda_2(T_{30})| = s^2$ admits a second, algebraic proof
from the CRT product structure, independent of explicit spectral computation.

\textbf{Route~2: CRT factorisation of eigenvalues.}
\label{thm:CRT}%
By CRT, $T_{30} \cong T_3 \otimes T_5$
in the Kronecker sense (the prime components are algebraically
independent by DPI, \cref{id:dpi}). The eigenvalues of a Kronecker product
are products of eigenvalues:
\[
  \mathrm{spec}(T_{30}) \supseteq
  \{\lambda_i(T_3) \cdot \lambda_j(T_5) : i,j\}.
\]
For $T_3 = \begin{pmatrix}0&1\\1&0\end{pmatrix}$:
$\mathrm{spec}(T_3) = \{+1, -1\}$.
For $T_5$ ($4\times 4$, row-stochastic): $\lambda_1(T_5) = 1$,
$|\lambda_2(T_5)| = 1/4$ (the T1 constraint---zero diagonal in $T_3$---propagates
via CRT to fix the spectral gap of $T_5$; without T1, the uniform
distribution would give $|\lambda_2| = 1/3$).
(Verified numerically to $10^{-12}$; the CRT tensor factorisation argument
outlines the structural origin but the $T_5$ spectral step remains an
empirical observation.)
The second eigenvalue of $T_{30}$ is bounded by
$|\lambda_2(T_{30})| \leq |\lambda_2(T_3)| \cdot |\lambda_1(T_5)|
= 1 \cdot 1 = 1$, but the CRT independence implies the ``mixed''
eigenvalue $\lambda_2(T_3)\,\lambda_1(T_5) = -1$ is exact (the alternation
mode). The next eigenvalue is $\lambda_1(T_3)\,\lambda_2(T_5) = 1/4 = s^2$.
This is the convergence-controlling eigenvalue: $|\lambda_2^{\rm eff}| = s^2$.

\textbf{Primorial scaling.}
The same CRT argument generalizes: for any primorial $m_k = \prod_{p \leq p_k} p$,
the convergence rate is $|\lambda_2^{\rm eff}(T_{m_k})| \leq s^2$, with
equality at $m = 30$ (the level where $T_5$ first contributes). This explains
why $s^2$ is a \emph{universal} spectral bound: it arises from the CRT tensor
structure, not from specific numerics of $T_{30}$.
\end{thmbox}

% ============================================================
\subsection{Vertex--edge bifurcation}
\label{sec:bifurcation}

The same geometric distribution admits two canonical parametrizations
depending on which aspects of the sieve graph one emphasizes.

\mTHM
\begin{theorem}[Vertex--Edge Bifurcation]
\label{thm:bifurcation}
At the fixed point $\mustar = 15$ (the companion article~\cite{companion_M3}), the sieve admits exactly
two canonical deformation parameters:

\textbf{Vertex branch (max-entropy):}
\begin{equation}
  \qstat = 1 - \frac{2}{\mustar} = \frac{13}{15} \approx 0.8\overline{6}.
  \label{eq:qstat_val}
\end{equation}

\textbf{Edge branch (Gibbs):}
\begin{equation}
  \qtherm = e^{-1/\mustar} = e^{-1/15} \approx 0.9355.
  \label{eq:qtherm_val}
\end{equation}

The bifurcation gap (latent heat) is:
\begin{equation}
  L = \qtherm - \qstat = e^{-1/15} - \frac{13}{15} \approx 0.069.
  \label{eq:latent_heat}
\end{equation}

\begin{proof}
\textbf{Vertex branch:} from L0 (Theorem~\ref{thm:L0}).

\textbf{Edge branch:} the Gibbs weight for a gap of size $g$ at
inverse temperature $\beta$ is $e^{-\beta g}$. The value
$\beta = 1/\mustar$ is not assumed but derived below from the
self-consistency condition: the unique $\beta$ for which the
geometric distribution is Gibbs-consistent is $\beta = 1/\mustar$.
For the geometric distribution with parameter $q$, the Gibbs-weighted
normalization condition is:
\[
  \mathbb{E}[e^{-g/\mustar}] = \frac{(1-q)e^{-1/\mustar}}{1-q\,e^{-1/\mustar}},
\]
which is self-consistent (equals $q$ in the edge formulation) when
$q = e^{-1/\mustar}$. This is the unique Gibbs-consistent parameter.

\textbf{Distinctness:} $\qtherm \neq \qstat$ iff $e^{-1/15} \neq 13/15$,
which is true (Taylor: $e^{-x} = 1 - x + x^2/2 - \ldots$; at $x = 1/15$:
$e^{-1/15} \approx 0.9355 > 0.8\overline{6} = \qstat$). \qed
\end{proof}
\end{theorem}

\begin{remark}[Derivation of the Gibbs parameter]
\label{rem:gibbs_derivation}
The Gibbs parameter $\beta = 1/\mustar$ is the unique inverse
temperature satisfying the thermodynamic consistency condition
$\partial F / \partial \beta |_{\beta = 1/\mustar} = 0$, where
$F = -\beta^{-1} \ln Z$ is the free energy of the transfer operator
$T_m$.  The full derivation is given in
Article~III~\cite{companion_M3}.
\end{remark}

\begin{remark}[Three independent routes to the bifurcation]
\label{rem:bifurcation_routes}
The existence of exactly two canonical parameters admits three
independent derivations, using disjoint mathematical machinery.

\textbf{Route~1: Optimisation / statistical (above).}
$\qstat$ from Lagrange multipliers (max-entropy under mean constraint);
$\qtherm$ from Gibbs self-consistency ($q = e^{-\beta}$ at inverse
temperature $\beta = 1/\mustar$).

\textbf{Route~2: Exponential family duality (information geometry).}
The geometric distribution $p_k = (1-q)\,q^{k-1}$ is an exponential
family with natural parameter $\theta = \ln q$ and sufficient statistic
$T(k) = k$. Every exponential family admits exactly two dual coordinate
systems (Amari~\cite{Amari2016}):
\begin{itemize}[nosep]
\item \emph{Expectation (moment) coordinate}:
  $\eta = \mathbb{E}[k] = 1/(1-q) = \mustar/2$.
  Solving gives $q = 1 - 2/\mustar = \qstat$.
\item \emph{Natural (canonical) coordinate}:
  $\theta = \ln q = -\beta$, with $\beta = 1/\mustar$.
  Exponentiating gives $q = e^{-1/\mustar} = \qtherm$.
\end{itemize}
The two parameters are related by the Legendre transform of the
log-partition function $\psi(\theta) = -\ln(1 - e^\theta)$:
\begin{equation}
  \eta = \psi'(\theta) = \frac{e^\theta}{1 - e^\theta}
       = \frac{1}{e^{-\theta} - 1}.
  \label{eq:legendre_duality}
\end{equation}
The bifurcation $\qstat \neq \qtherm$ is therefore equivalent to the
\emph{nonlinearity of the Legendre transform}: $\eta(\theta)$ is convex,
not affine. This is a structural necessity of any non-degenerate
exponential family---the two branches \emph{must} exist and \emph{must}
differ (Figure~\ref{fig:free_energy_3d}).

\textbf{Route~3: Partition function (statistical mechanics).}
The canonical partition function
$Z(\beta) = \sum_{k=1}^\infty e^{-\beta k} = e^{-\beta}/(1-e^{-\beta})$
defines the free energy $F(\beta) = -\ln Z(\beta)$. The mean is
$\langle k \rangle = -\partial \ln Z / \partial\beta = 1/(e^\beta - 1) + 1$.
Setting $\beta = 1/\mustar$ directly gives $q = e^{-\beta} = \qtherm$.
This is the standard canonical ensemble route, independent of both the
Jaynes max-entropy argument (Route~1) and the information-geometric
characterization (Route~2).

\medskip\noindent
\textbf{Summary.}
\begin{center}
\begin{tabular}{llll}
\toprule
\textbf{Route} & \textbf{Domain} & \textbf{$\qstat$ from\ldots} & \textbf{$\qtherm$ from\ldots} \\
\midrule
Optimisation & Calculus of variations & Lagrange (mean $= \mu/2$) & Gibbs self-consistency \\
Exp.\ family & Information geometry & Moment coordinate $\eta$ & Natural coordinate $\theta$ \\
Partition fn.\ & Statistical mechanics & $\langle k \rangle = \mu/2$ & $Z(\beta)$, $\beta = 1/\mu$ \\
\bottomrule
\end{tabular}
\end{center}
The three routes use disjoint mathematics and converge on the same
pair $(\qstat, \qtherm)$.
Route~2 also explains \emph{why exactly two} branches exist: any
one-parameter exponential family has exactly two dual coordinates.
\end{remark}

\begin{remark}[Physical meaning of the bifurcation]
The vertex/edge duality is the structural origin of the lepton/quark
distinction in the companion article~\cite{companion_M5}. The two branches correspond to
two dual readings of the same sieve graph:
\begin{center}
\begin{tabular}{lll}
\toprule
& \textbf{Vertex (nodes)} & \textbf{Edge (arcs)} \\
\midrule
Parameter $q$ & $\qstat = 13/15$ & $\qtherm = e^{-1/15}$ \\
Sieve object & Residue classes & Transition weights \\
Physics sector & Leptonic & Quark (BRIDGE) \\
\bottomrule
\end{tabular}
\end{center}
The physical identification is a bridge claim (the companion article~\cite{companion_M5}, \BRIDGETAG).
Here we establish only the existence of the two branches
(Figure~\ref{fig:bifurcation}).
\end{remark}

\begin{figure}[htbp]
\centering
\begin{tikzpicture}[>=stealth]
  % Rescaled coordinates: q mapped to x via x = (q - 0.78) * 55
  % q=0.80 -> x=1.1, q=0.85 -> x=3.85, q=0.8667 -> x=4.77
  % q=0.90 -> x=6.6, q=0.9355 -> x=8.55, q=0.95 -> x=9.35, q=1.00 -> x=12.1
  \pgfmathsetmacro{\K}{55}   % horizontal scale factor
  \pgfmathsetmacro{\qmin}{0.78}
  \newcommand{\qx}[1]{\pgfmathparse{(#1-\qmin)*\K}\pgfmathresult}
  % axes
  \draw[->,thick] (0,0) -- (13.2,0) node[right,font=\normalsize] {$q$};
  \draw[->,thick] (1.1,-0.4) -- (1.1,5.5);
  % q-axis ticks
  \foreach \q in {0.80, 0.85, 0.90, 0.95, 1.00} {
    \pgfmathsetmacro{\xx}{(\q-\qmin)*\K}
    \draw (\xx, -0.12) -- (\xx, 0.12);
    \node[below, font=\small] at (\xx, -0.18) {$\q$};
  }
  % q_stat = 13/15 = 0.8667
  \pgfmathsetmacro{\xstat}{(0.8667-\qmin)*\K}
  \draw[thmblue, very thick] (\xstat, -0.15) -- (\xstat, 3.8);
  \fill[thmblue] (\xstat, 2.2) circle (2.5pt);
  \node[above, font=\normalsize, thmblue, align=center] at (\xstat, 4.0)
    {\textbf{Vertex}\\[2pt]$\qstat = 13/15$};
  % q_therm = e^{-1/15} = 0.9355
  \pgfmathsetmacro{\xtherm}{(0.9355-\qmin)*\K}
  \draw[predred, very thick] (\xtherm, -0.15) -- (\xtherm, 3.8);
  \fill[predred] (\xtherm, 2.2) circle (2.5pt);
  \node[above, font=\normalsize, predred, align=center] at (\xtherm, 4.0)
    {\textbf{Edge}\\[2pt]$\qtherm = e^{-1/15}$};
  % Latent heat gap
  \draw[<->, condorange, very thick] (\xstat, 2.2) -- (\xtherm, 2.2);
  \pgfmathsetmacro{\xmid}{(\xstat+\xtherm)/2}
  \node[above, font=\normalsize, condorange] at (\xmid, 2.35)
    {$L = 0.069$};
  % Sector labels
  \pgfmathsetmacro{\xleft}{(\xstat-1.0)}
  \pgfmathsetmacro{\xright}{(\xtherm+1.2)}
  \node[font=\small, thmblue] at (\xleft, 1.3) {max-entropy};
  \node[font=\small, thmblue] at (\xleft, 0.8) {(nodes)};
  \node[font=\small, predred] at (\xright, 1.3) {Gibbs};
  \node[font=\small, predred] at (\xright, 0.8) {(arcs)};
  % Physical annotation (bridge)
  \node[font=\small, gray] at (\xleft, 0.2) {leptonic {\footnotesize(BRIDGE)}};
  \node[font=\small, gray] at (\xright, 0.2) {quark {\footnotesize(BRIDGE)}};
\end{tikzpicture}
\caption{Vertex--edge bifurcation at $\mustar = 15$.
  The same geometric distribution admits two canonical parameters:
  $\qstat = 13/15$ (vertex/max-entropy, blue) and
  $\qtherm = e^{-1/15}$ (edge/Gibbs, red).
  The latent heat $L = \qtherm - \qstat \approx 0.069$ (orange) seeds
  the structural distinction between two physical sectors.}
\label{fig:bifurcation}
\end{figure}

\begin{figure}[htbp]
\centering
\begin{tikzpicture}[xscale=1.6, yscale=1.1, >=stealth]
  % --- Three stacked cross-sections of F(q) at different mu ---
  % Axes
  \draw[->, thick] (-0.5,0) -- (7.5,0) node[right, font=\small] {$q$};
  \draw[->, thick] (0,-0.3) -- (0,7.5) node[above, font=\small] {$F(q,\mu)$};
  % --- Bottom: mu < mu* (single well, gray) ---
  \draw[gray!60, thick, domain=0.5:6.5, samples=60, smooth]
    plot (\x, {0.4 + 0.15*(\x-3.5)*(\x-3.5)});
  \node[right, font=\scriptsize, gray] at (6.8, 1.8)
    {$\mu < \mustar$};
  % --- Middle: mu = mu* (double well, orange) — shifted up ---
  \draw[condorange, very thick, domain=0.5:6.5, samples=60, smooth]
    plot (\x, {3.5 - 1.0*cos((\x-3.5)*50)});
  % Two minima at mu*
  \fill[thmblue] (2.5, 2.55) circle (2pt);
  \node[below left, font=\scriptsize, thmblue] at (2.2, 2.4) {$\qstat$};
  \fill[predred] (4.5, 2.55) circle (2pt);
  \node[below right, font=\scriptsize, predred] at (4.8, 2.4) {$\qtherm$};
  % Barrier marker
  \fill[condorange] (3.5, 4.45) circle (1.5pt);
  \node[above, font=\scriptsize, condorange] at (3.5, 4.6) {barrier};
  % Latent heat arrow
  \draw[<->, condorange, semithick] (2.5, 2.1) -- (4.5, 2.1);
  \node[below, font=\scriptsize, condorange] at (3.5, 1.9) {$L = 0.069$};
  \node[right, font=\scriptsize, condorange] at (6.8, 3.5)
    {$\mu = \mustar$};
  % --- Top: mu >> mu* (deep double well, blue) — shifted up more ---
  \draw[thmblue, thick, domain=0.5:6.5, samples=60, smooth]
    plot (\x, {6.0 - 1.5*cos((\x-3.5)*50)});
  \node[right, font=\scriptsize, thmblue] at (6.8, 6.0)
    {$\mu > \mustar$};
  % Vertical arrows showing evolution
  \draw[->, gray, thick, densely dotted] (7.2, 1.5) -- (7.2, 3.2);
  \draw[->, gray, thick, densely dotted] (7.2, 4.0) -- (7.2, 5.5);
  \node[right, font=\scriptsize, gray, rotate=90] at (7.4, 3.5) {increasing $\mu$};
  % q-axis ticks
  \node[below, font=\tiny] at (2.5, -0.05) {$0.87$};
  \node[below, font=\tiny] at (4.5, -0.05) {$0.94$};
\end{tikzpicture}
\caption{Free energy cross-sections $F(q)$ at three values of $\mu$ (schematic).
  \textbf{Bottom} (gray): $\mu < \mustar$, single minimum---no phase distinction.
  \textbf{Middle} (orange): $\mu = \mustar$, two minima appear at
  $\qstat = 13/15 = 0.867$ (blue, vertex sector) and
  $\qtherm = e^{-1/15} = 0.936$ (red, edge sector),
  separated by a barrier with latent heat $L = 0.069$.
  \textbf{Top} (blue): $\mu > \mustar$, the double well deepens.
  This first-order bifurcation seeds the two coupling
  constants $\alphaEM$ and $\alpha_s$.}
\label{fig:free_energy_3d}
\end{figure}

% ============================================================
\subsection{Numerical values}
\label{sec:ch3_numerics}

\begin{center}
\begin{tabular}{lll}
\toprule
\textbf{Quantity} & \textbf{Value} & \textbf{Source} \\
\midrule
$\qstat$ & $13/15 = 0.86\overline{6}$ & L0 at $\mustar = 15$ \\
$\qtherm$ & $e^{-1/15} = 0.93562\ldots$ & Gibbs at $\mustar = 15$ \\
$L = \qtherm - \qstat$ & $0.06896\ldots$ & Bifurcation gap \\
$\alpha_{\rm cons}$ & $1/4 = s^2$ & T2 (spectral of $T_{30}$) \\
$\eta$ (efficiency) & $(\mustar - 1)/\mustar = 14/15 = 93.3\%$ & Carnot analog \\
$C_v$ & $N_{\rm gen} \cdot s = 3 \cdot 1/2 = 3/2$ & Heat capacity (BRIDGE) \\
\bottomrule
\end{tabular}
\end{center}

% ============================================================
\subsection{Résumé et connexions}
\label{sec:ch3_summary}

\begin{ledgerbox}
Hard core:
  L0 (geometric = unique max-entropy, standard Lagrange argument).
  Conservation identity (telescoping, exact).
  Bifurcation ($q_{\rm stat} = 13/15$ vertex; $q_{\rm therm} = e^{-1/15}$ edge).

Numerical (algebraic proof open):
  T2 ($\alpha_{\rm cons} = s^2 = 1/4$, spectral computation on $T_{30}$,
  verified to $10^{-12}$).

Conditional:
  The substitution $\mustar = 15$ in L0 uses the companion article~\cite{companion_M3} (T5), which is
  proved by scan (THM by scan, not analytically).

Bridge claim:
  Identification vertex $\leftrightarrow$ leptons, edge $\leftrightarrow$ quarks (the companion article~\cite{companion_M5}).
\end{ledgerbox}

\medskip\noindent\textbf{Connexions prospectives.}
\begin{itemize}
\item \cref{sec:gft} (GFT) assembles $\qstat$ and $\DKL$ into the Gap
  Fundamental Theorem: $\log_2 m = \DKL + H$.
\item The companion article~\cite{companion_M2} (holonomy) uses $\qstat$ and $\qtherm$ to define
  $\sin^2(\theta_p)$ at each prime.
\item The companion article~\cite{companion_M5} (couplings) uses the bifurcation to derive
  $\alphas$ from the edge branch.
\end{itemize}


% ============================================================
%  S5. The Gap Fundamental Theorem
% ============================================================
\section{Le Théorème Fondamental des Lacunes}
\label{sec:gft}

\epigraph{Information is physical.}{Rolf Landauer}

% ---- Status Box (R1) ----------------------------------------
\begin{statusbox}
\textbf{GOAL:}
  Prove the Gap Fundamental Theorem (GFT: $\log_2 m = \DKL + H$)
  as an exact algebraic identity and derive its consequences
  (Ruelle equivalence, Arrow of Time, Bekenstein bound).

\textbf{INPUTS:}
  \cref{sec:sieve} (transfer matrix $T_m$, DPI). \cref{sec:conservation} (conservation identity).

\textbf{MAIN RESULTS:}
  GFT-ID: $\log_2 m = \DKL(P\|U_m) + H(P)$ (exact identity, residual $<10^{-15}$).
  GFT*: structural stability of the partition (COND on PNT).
  Ruelle: $Z_{\mathrm{Ruelle}} = \mathrm{Tr}(T_m^N)$, $F_R = 0$ (IDENTITY).
  Arrow of Time: $dH/d\DKL = -1$ (IDENTITY).
  Bekenstein: $\DKL \leq \log_2 m$ (THM).

\textbf{STATUS:}
  \IDTAG~(GFT-ID, Arrow) + \CONDTAG~(Ruelle: conditional on irreducibility) + \THMTAG~(Bekenstein); \CONDTAG~(GFT*)

\textbf{NOT PROVED:}
  GFT-ID is a tautology of information theory (true for \emph{any} distribution).
  The depth of the identity is in its application, not the algebra.
  Physical thermodynamic interpretation (the companion article~\cite{companion_M5}).
\end{statusbox}

% ============================================================
\subsection{Intuition: the information budget}
\label{sec:ch4_intuition}

The sieve of Eratosthenes distributes primes non-uniformly across
residue classes. If we measure the total information capacity at
modulus $m$ as $\log_2 m$ (the entropy of the uniform distribution),
the GFT says: this capacity is spent in exactly two ways.

\begin{enumerate}
\item \textbf{Structure} $\DKL$: how far the actual prime distribution
  deviates from uniform. High $\DKL$ means the sieve has imposed strong
  constraints.
\item \textbf{Disorder} $H$: the residual uncertainty in the distribution.
  High $H$ means the distribution is spread out.
\end{enumerate}

The GFT says $\DKL + H = \log_2 m$ exactly, with zero residual. This
is not deep mathematics---it is a three-line tautology of information
theory. The content arrives when $P$ is the prime residue distribution:
then $\DKL$ and $H$ acquire the roles of \emph{sieve order} and
\emph{sieve entropy}, and the identity becomes the first law of
sieve thermodynamics.

\paragraph{Warning on depth.}
The GFT-ID is not a remarkable theorem about primes. It holds for
\emph{any} probability distribution on $\{0, \ldots, m-1\}$. The
claim ``GFT is deep'' is a claim about the \emph{interpretation} of
$\DKL$ and $H$ in the sieve context, not about the algebra.
The companion article~\cite{companion_M5} develops the thermodynamic reading; this section proves only
the identity and its immediate corollaries.

% ============================================================
\subsection{Statement and proof}
\label{sec:gft_statement}

\begin{definition}[Sieve distribution at modulus $m$]
For a prime $p$ chosen uniformly at random from the first $N$ primes,
the \emph{sieve distribution} at square-free modulus $m$ is
$P_m = (p_0, \ldots, p_{m-1})$ where
$p_r = \#\{p_k \leq N : p_k \equiv r \pmod m\}/N$.
The uniform reference is $U_m = (1/m, \ldots, 1/m)$.
\end{definition}

\mID
\begin{idbox}[Gap Fundamental Theorem (GFT-ID)]
\label{id:gft}
For \emph{any} probability distribution $P = (p_0, \ldots, p_{m-1})$
on $m$ classes with uniform reference $U_m$:
\begin{equation}
  \log_2 m = \DKL(P \| U_m) + H(P),
  \label{eq:gft}
\end{equation}
where $\DKL(P\|U_m) = \sum_r p_r \log_2(m\, p_r)$ and
$H(P) = -\sum_r p_r \log_2 p_r$.

\begin{proof}
Expand $\DKL$:
\begin{align*}
  \DKL(P\|U_m)
  &= \sum_r p_r \log_2\frac{p_r}{1/m}
   = \sum_r p_r \bigl(\log_2 m + \log_2 p_r\bigr) \\
  &= \log_2 m \cdot \underbrace{\sum_r p_r}_{=1}
   + \underbrace{\sum_r p_r \log_2 p_r}_{=-H(P)} \\
  &= \log_2 m - H(P).
\end{align*}
Hence $\DKL + H = \log_2 m$. The residual is zero by algebra; numerical
verification at $m = 210$ gives $|\text{residual}| < 10^{-15}$ bits.
\end{proof}
\end{idbox}

\begin{remark}[Universality]
The GFT holds for \emph{any} distribution $P$: Gaussian, Poisson,
random, or prime. The word ``Fundamental'' refers to its role in PT,
not to any claim of mathematical novelty. The identity is Shannon's
decomposition theorem~\cite{Shannon1948}.
\end{remark}

\paragraph{Numerical verification at primorial levels.}
Figure~\ref{fig:gft_budget} visualises the information budget at the
first four primorial levels.

\begin{center}
\begin{tabular}{ccccc}
\toprule
$m$ & $\log_2 m$ & $\DKL$ & $H$ & Residual \\
\midrule
$2$   & 1.000 & 1.000 & 0.000 & $< 10^{-15}$ \\
$6$   & 2.585 & 1.585 & 1.000 & $< 10^{-15}$ \\
$30$  & 4.907 & 3.252 & 1.655 & $< 10^{-15}$ \\
$210$ & 7.714 & 4.916 & 2.798 & $< 10^{-15}$ \\
\bottomrule
\end{tabular}
\end{center}

\begin{figure}[htbp]
\centering
\begin{tikzpicture}[xscale=1.8, yscale=0.85, >=stealth]
  % axes
  \draw[->,thick] (0,0) -- (5.5,0) node[right,font=\small] {primorial $m$};
  \draw[->,thick] (0,0) -- (0,8.5) node[above,font=\small] {bits};
  % y-axis ticks
  \foreach \y in {1,...,8} {
    \draw (-0.08, \y) -- (0.08, \y);
    \node[left, font=\tiny] at (-0.1, \y) {$\y$};
  }
  % bar width
  \pgfmathsetmacro{\bw}{0.5}
  % m=2: log2=1.000, DKL=1.000, H=0.000
  \fill[predred!70] (1-\bw/2, 0) rectangle (1+\bw/2, 1.000);
  \node[font=\tiny, white] at (1, 0.5) {$D_{\rm KL}$};
  \node[above, font=\scriptsize] at (1, 1.05) {$1.00$};
  \node[below, font=\scriptsize] at (1, -0.15) {$2$};
  % m=6: log2=2.585, DKL=1.585, H=1.000
  \fill[predred!70] (2-\bw/2, 0) rectangle (2+\bw/2, 1.585);
  \fill[thmblue!60] (2-\bw/2, 1.585) rectangle (2+\bw/2, 2.585);
  \node[font=\tiny, white] at (2, 0.8) {$D_{\rm KL}$};
  \node[font=\tiny, white] at (2, 2.1) {$H$};
  \node[above, font=\scriptsize] at (2, 2.65) {$2.59$};
  \node[below, font=\scriptsize] at (2, -0.15) {$6$};
  % m=30: log2=4.907, DKL=3.252, H=1.655
  \fill[predred!70] (3-\bw/2, 0) rectangle (3+\bw/2, 3.252);
  \fill[thmblue!60] (3-\bw/2, 3.252) rectangle (3+\bw/2, 4.907);
  \node[font=\tiny, white] at (3, 1.6) {$D_{\rm KL}$};
  \node[font=\tiny, white] at (3, 4.1) {$H$};
  \node[above, font=\scriptsize] at (3, 4.97) {$4.91$};
  \node[below, font=\scriptsize] at (3, -0.15) {$30$};
  % m=210: log2=7.714, DKL=4.916, H=2.798
  \fill[predred!70] (4-\bw/2, 0) rectangle (4+\bw/2, 4.916);
  \fill[thmblue!60] (4-\bw/2, 4.916) rectangle (4+\bw/2, 7.714);
  \node[font=\tiny, white] at (4, 2.5) {$D_{\rm KL}$};
  \node[font=\tiny, white] at (4, 6.3) {$H$};
  \node[above, font=\scriptsize] at (4, 7.78) {$7.71$};
  \node[below, font=\scriptsize] at (4, -0.15) {$210$};
  % legend
  \fill[predred!70] (4.6, 7.5) rectangle (4.9, 7.8);
  \node[right, font=\scriptsize] at (4.95, 7.65) {$D_{\rm KL}$ (structure)};
  \fill[thmblue!60] (4.6, 6.8) rectangle (4.9, 7.1);
  \node[right, font=\scriptsize] at (4.95, 6.95) {$H$ (disorder)};
  % GFT identity annotation
  \node[font=\footnotesize, align=center] at (3, -1.0)
    {$\log_2 m = D_{\rm KL}(P\|U_m) + H(P)$\quad (residual $< 10^{-15}$)};
\end{tikzpicture}
\caption{Information budget of the sieve at primorial levels.
  The total capacity $\log_2 m$ (bar height) is exactly partitioned
  into structure ($D_{\rm KL}$, red) and disorder ($H$, blue) by the
  GFT identity. As $m$ grows, $D_{\rm KL}$ dominates: the sieve
  imposes increasing order.}
\label{fig:gft_budget}
\end{figure}

% ============================================================
\subsection{Structural stability: GFT*}
\label{sec:gft_star}

The GFT is universal, but the prime sieve has a further property:
the partition ratio $\DKL/H$ converges.

\mDER
\begin{theorem}[Structural Stability, GFT* (conditional on PNT)]
\label{thm:gft_star}
For the prime gap distribution at primorial levels $m_k = 2 \cdot 3 \cdots p_k$,
the partition ratio $\DKL(P_{m_k})/H(P_{m_k})$ converges to a definite
limit as $k \to \infty$. The partition is structurally stable.

\begin{proof}
By the prime number theorem for arithmetic progressions (PNT-AP), the
distribution $P_{m_k}$ of primes modulo $m_k$ converges to the Euler
product measure $\mu_{m_k}(r) \propto \prod_{p|m_k}(1-1/p)^{-1}$
(for $\gcd(r, m_k) = 1$), which is uniform on reduced residues.
The KL divergence $\DKL(P_{m_k}\|\mu_{m_k})$ and entropy $H(P_{m_k})$
are continuous functions of $P_{m_k}$ and converge as $k \to \infty$.
Stability follows from the strict convexity of $\DKL$ and concavity of $H$.
The convergence rate is $O(1/\ln p_k)$ by PNT-AP; the ratio
$\DKL/H \to L$ where $L \approx s^2 = 1/4$.
\end{proof}
\end{theorem}

\begin{remark}
GFT* distinguishes the prime sieve from a random sieve. A random process
satisfies GFT-ID (universally), but the ratio $\DKL/H$ fluctuates without
converging. For primes, it converges to the Mertens constant.
\end{remark}

\paragraph{Complex extension.}
The action $S_{\mathrm{PT}} = -\sum\ln\sinp{p}$ extends to
$S_{\mathbb{C}} = -\sum\ln w_p$ with $\mathrm{Re}(S_{\mathbb{C}}) = S_{\mathrm{PT}}/2$,
where $w_p = (1 - e^{2i\theta_p})/2$ is the complex sieve variable.
See the companion article~\cite{companion_M4}.

% ============================================================
\subsection{Ruelle equivalence}
\label{sec:ruelle}

The transfer matrix $T_m$ connects GFT to Ruelle's thermodynamic
formalism~\cite{Ruelle1978}.

\mID
\begin{idbox}[Ruelle Equivalence]
\label{id:ruelle}
The Ruelle partition function of the sieve is:
\begin{equation}
  Z_{\mathrm{Ruelle}} = \mathrm{Tr}(T_m^N) = \sum_i \lambda_i^N,
  \label{eq:ruelle}
\end{equation}
where $\lambda_i$ are the eigenvalues of $T_m$. For a row-stochastic
matrix, $\lambda_0 = 1$ (Perron--Frobenius). The Ruelle free energy is:
\begin{equation}
  F_{\mathrm{Ruelle}} = -\ln \lambda_0 = 0.
  \label{eq:ruelle_free}
\end{equation}
The system is in equilibrium: $F_R = 0$.

\medskip\noindent\textbf{Status:} \CONDTAG\ (conditional on irreducibility of $T_m$
for all $m$; verified numerically for $m \leq 2310$, $N \leq 10^7$,
but not proved for all $m$).

\begin{proof}
$T_m$ restricted to the $\varphi(m)$ coprime residue classes is a
row-stochastic matrix (rows sum to~1, entries $\geq 0$).
It is empirically irreducible for all moduli tested ($m \leq 2310$,
primes up to $10^7$): every entry $(T_m)_{ab}$ with $\gcd(a,m) = \gcd(b,m) = 1$
is strictly positive.
(A proof of irreducibility for all sample sizes would require showing
that every coprime-class pair is realized by consecutive primes, which
is widely expected but not formally established.)
Assuming irreducibility, the Perron--Frobenius theorem guarantees a
unique maximal eigenvalue $\lambda_0 = 1$ with positive eigenvector
(stationary distribution $\pi$).
The trace formula $\mathrm{Tr}(T_m^N) = \sum_i
\lambda_i^N$ is the spectral decomposition of the trace.
$F_R = -\ln 1 = 0$ follows directly.
\end{proof}
\end{idbox}

\begin{remark}
The same partition function $Z = \mathrm{Tr}(T_m^N)$ appears in
Polyakov's transfer-matrix method~\cite{Polyakov1981}
(1D statistical mechanics) and in Regge calculus~\cite{Regge1959}. The triple coincidence $Z_{\mathrm{Feynman}} =
Z_{\mathrm{Ruelle}} = Z_{\mathrm{Polyakov}}$ is developed in
the companion article~\cite{companion_M5}.
\end{remark}

\begin{remark}[Topological pressure route: second proof of $F_R = 0$]
\label{rem:ruelle_pressure}
The Ruelle free energy $F_R = 0$ admits a second, structurally
independent derivation via the \emph{topological pressure} of the
Ruelle--Perron--Frobenius formalism~\cite{Ruelle1978}.

Define the potential $\varphi : \Sigma \to \mathbb{R}$ on the symbolic
shift $\Sigma$ induced by $T_m$ as $\varphi(\omega) = \ln T_m(\omega_0, \omega_1)$
(the log-transition weight).  The topological pressure is:
\[
  P(\varphi) = \lim_{n \to \infty} \frac{1}{n}
    \ln \sum_{\omega \in \Sigma_n} \exp\!\Bigl(\sum_{k=0}^{n-1}\varphi(\sigma^k\omega)\Bigr)
  = \ln \lambda_0(L_\varphi),
\]
where $L_\varphi$ is the Ruelle transfer operator and $\lambda_0(L_\varphi)$
its spectral radius.  For a row-stochastic matrix, $\lambda_0 = 1$
(Perron--Frobenius), hence $P(\varphi) = \ln 1 = 0$.
The variational principle then gives:
\[
  0 = P(\varphi) = \sup_\mu \bigl[h(\mu) + \int \varphi\,d\mu\bigr],
\]
where $h(\mu)$ is the measure-theoretic entropy.
Under the irreducibility assumption (empirically verified, see proof above),
the supremum is attained at the unique equilibrium state $\mu_{\rm eq} = \pi$
(the Perron eigenvector).
Note that $\pi$ is the measure of maximal entropy (MME) if and only if
$T_m$ is doubly stochastic; this holds at the sieve level for $m = 3$
(where $T_3$ is the antidiagonal permutation matrix) but need not hold
for the empirical prime transfer matrix at general~$m$.
The equality $F_R = -P(\varphi) = 0$
recovers Identity~\ref{id:ruelle} without using the trace formula
$\mathrm{Tr}(T^N)$---the argument depends only on the spectral
radius of $L_\varphi$, not on the matrix decomposition.

This provides a second independent route to $F_R = 0$: the first
uses the trace and Perron--Frobenius eigenvalue directly; the second
uses the variational principle of topological pressure.
Both routes are conditional on irreducibility of the restricted
matrix $T_m$ on coprime classes.
\end{remark}

% ============================================================
\subsection{Arrow of Time}
\label{sec:arrow}

Differentiating the GFT at fixed $m$ gives an irreversibility statement.

\mID
\begin{idbox}[Arrow of Time]
\label{id:arrow}
At fixed modulus $m$, along any path in distribution space satisfying
the GFT constraint:
\begin{equation}
  \frac{dH}{d\DKL} = -1.
  \label{eq:arrow}
\end{equation}

\begin{proof}
Differentiate $\DKL + H = \log_2 m$ (constant in $m$) with respect
to any parameter: $d\DKL + dH = 0$, hence $dH/d\DKL = -1$.

More precisely: as the prime distribution evolves toward the uniform
(as $N \to \infty$ by PNT-AP), $\DKL$ decreases while $H$ increases,
at equal and opposite rates along the GFT line.
\end{proof}
\end{idbox}

\begin{remark}
The Arrow of Time (dH/dDKL = -1) is exact: no inequality, no approximation.
The sieve's second law is an identity rather than an inequality.
\end{remark}

% ============================================================
\subsection{Bekenstein bound}
\label{sec:bekenstein}

\mTHM
\begin{theorem}[Bekenstein Bound]
\label{thm:bekenstein}
For any distribution $P$ on $m$ classes:
\begin{equation}
  \DKL(P\|U_m) \leq \log_2 m,
  \label{eq:bekenstein}
\end{equation}
with equality iff $P$ is a delta function on one class.

\begin{proof}
Since $H(P) \geq 0$ (entropy non-negative, with equality iff $P$ is
a point mass), the GFT gives $\DKL = \log_2 m - H \leq \log_2 m$.
Equality: $H = 0$ iff $P = \delta_r$ for some $r$.
\end{proof}
\end{theorem}

% ============================================================
\subsection{The derivation chain}
\label{sec:ch4_chain}

\[
  s = \tfrac{1}{2}
  \xrightarrow{\text{T1}} T_3
  \xrightarrow{\text{L0}} P_m \text{ geometric}
  \xrightarrow{\text{GFT}} \log_2 m = \DKL + H \xrightarrow{\text{Arrow}} dH/d\DKL = -1
\]

This chain is entirely within arithmetic. No physics enters.

% ============================================================
\subsection{Zero as cycle closure}
\label{sec:ch4_zero}

The GFT identity $\log_2 m = \DKL + H$ has a natural interpretation
through the concept of \textbf{cycle closure}.

\begin{remark}[Zero is not absence but phase cancellation]
\label{rem:zero_phase_gft}
In the sieve, the residue $n \equiv 0 \pmod{p}$ does not mean
``$n$ is absent at level~$p$'': it means the $p$-phase of~$n$
has \textbf{completed a full cycle} of $\mathbb{Z}/p\mathbb{Z}$
and returned to the origin.  A composite $n = k \cdot p$ has
traversed the circle $\mathbb{Z}/p\mathbb{Z}$ exactly $k$~times;
its holonomy has cancelled.

This gives a two-level reading of the sieve:
\begin{itemize}[nosep]
\item \textbf{Primes} $=$ perpetually open phases (no cycle closes
  below them).  They are irreducible precisely because they carry
  \emph{unclosed} holonomy.
\item \textbf{Composites} $=$ partially closed phases.
  $n = \prod p_i^{a_i}$ means the cycles of $\{p_i\}$ have each
  closed $a_i$ times.  The sieve removes them because a closed
  cycle carries no further open information.
\item \textbf{Zero} $=$ universal closure: every prime divides~$0$,
  so every cycle is simultaneously closed.  It is the
  \emph{informational vacuum}---not empty, but fully cancelled.
\item \textbf{One} $=$ no closure: no prime divides~$1$, so every
  cycle is fully open.  It is the \emph{identity operator}.
\end{itemize}

The GFT identity reflects this structure.  When $\DKL = 0$
(complete dissipation), all structure has been returned to
entropy: every informational cycle has closed.  When
$H = 0$ (zero entropy), all entropy has been converted to
structure: no cycle has closed, the state is maximally
persistent.  The conservation law $\DKL + H = \log_2 m$ says
that the total capacity is shared between open phases (structure)
and closed phases (entropy), with no loss.
\end{remark}

\begin{remark}[Fixed points as global cycle closures]
\label{rem:global_closure}
The fixed-point equation $\mustar = \sum_{\mathrm{active}} p$
(the companion article~\cite{companion_M3}) is a \textbf{global} cycle closure:
the sum of the active operators equals the scale that defines them.
The sieve ``recognises itself''---the structure it generates
(the active primes) sums to the scale that generates the structure.

Local closure: $n \equiv 0 \pmod{p}$ (one prime divides $n$).\\
Global closure: $\mustar = \sum p_{\rm active}$ (the sieve closes on itself).
\end{remark}

\begin{remark}[GFT as addition/multiplication conservation]
\label{rem:gft_add_mult}
The GFT identity $\log_2 m = \DKL + H$ can be read as a
\textbf{conservation law between the two operations}:
\begin{itemize}[nosep]
\item $\DKL$ (structure, persistence) measures
  \emph{multiplicative} information: how far the distribution
  deviates from uniform, quantified via log-ratios (logarithm
  converts $\times$ to $+$).
\item $H$ (entropy, disorder) measures \emph{additive}
  information: the sum $-\sum P \log P$ counts bits
  additively.
\item $\log_2 m$ (capacity) is the total: fixed, independent
  of how the information is partitioned.
\end{itemize}
The arrow identity $dH/d\DKL = -1$ says:
every unit of multiplicative structure gained costs exactly one
unit of additive entropy.  The two operations are in exact
balance.

This is the information-theoretic expression of the duality
described in \cref{rem:add_mult_duality}: the sieve
alternates $\times$ and $+$, and GFT conserves their sum.
\end{remark}

% ============================================================
\subsection{Résumé et connexions}
\label{sec:ch4_summary}

\begin{ledgerbox}
Hard core:
  GFT-ID ($\log_2 m = D_{\rm KL} + H$) -- exact algebraic tautology.
  Arrow ($dH/dD_{\rm KL} = -1$) -- exact differentiation.
  Bekenstein ($D_{\rm KL} \leq \log_2 m$) -- follows from $H \geq 0$.
  GFT* (structural stability) -- uses PNT-AP convergence.

Conditional:
  Ruelle ($F_R = 0$, $\mathrm{Tr}(T^N)$) -- conditional on irreducibility
  of $T_m$ for all $m$; verified numerically for $m \leq 2310$.

Bridge claim:
  Thermodynamic interpretation of $D_{\rm KL}$ and $H$ (the companion article~\cite{companion_M5}).
\end{ledgerbox}

\medskip\noindent\textbf{Connexions prospectives.}
\begin{itemize}
\item The companion article~\cite{companion_M2} identifies $\DKL$ as the canonical divergence (Shore--Johnson).
\item The companion article~\cite{companion_M5} uses $Z_{\mathrm{Ruelle}} = \mathrm{Tr}(T_m^N)$ to
  construct the quantum partition function.
\item The companion article~\cite{companion_M5} promotes GFT to the first law of sieve thermodynamics.
\end{itemize}


% ============================================================
%  Conclusion
% ============================================================
\section{Conclusion}
\label{sec:conclusion}

This article has established the first four pillars of the Theory of
Persistence, proceeding from the single observation that the
Eratosthenes sieve forbids self-transitions modulo~3 among its
candidates.

\textbf{The sieve as a dynamical field} (Section~\ref{sec:sieve}).
The gap sequence $\{g_n\}$ is the unique degree of freedom.  Its
mod-$m$ projection defines a $\mathbb{Z}/m\mathbb{Z}$-valued gauge
connection, factoring via CRT into independent components at each
prime.  The forbidden-transition theorem (T1) forces the mod-$3$
transfer matrix to be purely antidiagonal, $T_3 = \mathrm{antidiag}(1,1)$,
with eigenvalues $\{+1,-1\}$ and stationary distribution $s = 1/2$.
Three independent routes---gap arithmetic, $\mathbb{Z}/6\mathbb{Z}$
involution, spectral constraint---establish $T_3$; the DPI
monotonicity orders the sieve levels by KL divergence.

\textbf{Uniqueness of Eratosthenes} (Section~\ref{sec:uniqueness}).
Four naturality theorems (N1--N4) prove that Eratosthenes is the
unique self-consistent, structurally minimal sieve with canonical
ordering; $p = 3$ is forced as the first cascade level and $s = 1/2$
as the first dynamical output.  Theorem T6 then provides three
logically independent proofs of sieve uniqueness---from field
structure (T6a: the only proper ideal of $\mathbb{Z}/p\mathbb{Z}$ is
$\{0\}$), from four ring-congruence axioms C1--C4 (T6b: independence
verified by 4/4 counter-models), and from canonical information
geometry (T6c: $\DKL$ unique by Shore--Johnson, Fisher metric unique
by \v{C}encov).

\textbf{Conservation laws} (Section~\ref{sec:conservation}).
The maximum-entropy gap distribution is geometric (L0, by strict
concavity of entropy under a mean constraint).  The telescoping
conservation identity $\sum g_n = p_{N+1} - 2$ is the sieve's equation
of state.  The spectral conservation exponent $\alpha_{\rm cons} =
s^2 = 1/4$ (Theorem T2) controls the exponential convergence rate,
established both by direct spectral computation on $T_{30}$ and by a
CRT trace formula.  The vertex--edge bifurcation at $\mustar = 15$
produces exactly two canonical parameters, $\qstat = 13/15$ and
$\qtherm = e^{-1/15}$, derived independently via three routes
(optimisation, exponential-family duality, partition function).

\textbf{The Gap Fundamental Theorem} (Section~\ref{sec:gft}).
The exact algebraic identity $\log_2 m = \DKL + H$ decomposes the
total information capacity into persistent structure and irreducible
noise.  The Arrow of Time $dH/d\DKL = -1$ is an exact identity, not
an inequality.  The Ruelle identification connects the GFT to
thermodynamic formalism: the sieve's transfer matrix has zero free
energy ($F_R = 0$), and its stationary distribution is the unique
equilibrium measure.  The Bekenstein bound $\DKL \leq \log_2 m$
follows as a corollary.

\medskip
\noindent\textbf{Open problems.}
\begin{enumerate}[nosep]
\item \emph{Convergence rate of $\alpha_k$.}  The value $s = 1/2$ is
  the unique stationary distribution of $T_3$ (a theorem of linear
  algebra).  The rate at which the empirical frequency $\alpha_k$
  converges to $s$ is quantified in Article~III~\cite{companion_M3}
  via spectral annihilation and Mertens compactness.
\item \emph{Irreducibility of $T_m$.}  The Ruelle identification and
  GFT* both assume irreducibility of the transfer matrix on coprime
  classes.  A proof for all $m$ remains open.
\item \emph{Fixed point $\mustar = 15$.}  The self-consistency scan
  identifies $\mustar = 15$ as the unique cascade fixed point; an
  analytic proof is developed in the companion article~\cite{companion_M3}.
\item \emph{Physical interpretation.}  All results in this article are
  purely arithmetic.  Whether the vertex--edge bifurcation, the
  conservation laws, and the GFT identity have physical content is
  the subject of the companion article~\cite{companion_M5}.
\end{enumerate}

\medskip
\noindent\textbf{What comes next.}
The companion article~\cite{companion_M2} develops the canonical information geometry
(Fisher metric, holonomy angles $\sin^2\theta_p$) and the internal
angle scale.  The companion article~\cite{companion_M3} establishes the exact recurrences,
spectral annihilation, and the fixed-point theorem $\mustar = 15$.
Together, the five articles provide the complete mathematical
foundation.  The bridge to physical observables (43 Standard Model
quantities at $0.30\%$ mean deviation, zero fitted parameters) is
developed in Article~V~\cite{companion_M5}.

\begin{table}[htbp]
\centering
\caption{Dependency and epistemic status of main results.}
\small
\begin{tabular}{@{}llll@{}}
\toprule
\textbf{Result} & \textbf{Depends on} & \textbf{Status} & \textbf{Notes} \\
\midrule
T1 (Forbidden Trans.) & None & \THMTAG & Unconditional \\
T3 (Antidiagonal) & T1 & \THMTAG & Unconditional \\
$s = 1/2$ & T3 (Perron--Frobenius) & \THMTAG & Unique stationary distribution \\
N1--N2 (Uniqueness) & FTA & \THMTAG & Classical \\
N3--N4 (Minimality) & N1--N2 & \THMTAG & Novel \\
T6 (Sieve uniqueness) & N1--N4 & \THMTAG & Three proofs \\
L0 (Max entropy) & Lagrange, Jaynes & \THMTAG & Standard technique \\
T2 (Spectral cons.) & $T_{30}$ spectrum & \VALTAG & Numerical; alg.\ proof open \\
Bifurcation & L0 + $\mustar$ (M3) & \CONDTAG & Conditional on M3 \\
GFT identity & Shannon decomp. & \IDTAG & Valid for any distribution \\
Ruelle equivalence & Irreducibility & \CONDTAG & Open for all $m$ \\
\bottomrule
\end{tabular}
\end{table}


\bibliographystyle{plainnat}
\bibliography{references}

\end{document}
